% ---------------------------------------------------------------------------
% Chapter 8: The Crowd on the Road: Pileup
% The scripts, records, logs and films of this stop live in example/PU/, and
% the technical companion of the chapter is chapter/chapter_08_PU.md: how to
% run the scripts, how to read the records and logs, how the official numbers
% were obtained, and the full tables. The book gives the physics.
% The crowd is the crowd of Run 3 (2024): sixty collisions per crossing at
% 13.6 TeV, the 2024 beam spot, the production window bx = -12 ... +3.
% ---------------------------------------------------------------------------
%% chapter-local shorthands
\providecommand{\akfour}{\code{ak4GenJets}\xspace}
\providecommand{\akfournonu}{\code{ak4GenJetsNoNu}\xspace}
\providecommand{\mup}{\ensuremath{\mu}\xspace}
\providecommand{\mutrue}{\ensuremath{\mu_{\mathrm{true}}}\xspace}
\providecommand{\rhoPU}{\ensuremath{\rho}\xspace}
\providecommand{\sigmainel}{\ensuremath{\sigma_{\mathrm{inel}}}\xspace}
\providecommand{\bx}{\ensuremath{\mathrm{bx}}\xspace}
\providecommand{\ns}{\ensuremath{\,\mathrm{ns}}\xspace}
\providecommand{\ps}{\ensuremath{\,\mathrm{ps}}\xspace}
\providecommand{\cm}{\ensuremath{\,\mathrm{cm}}\xspace}
\providecommand{\um}{\ensuremath{\,\mu\mathrm{m}}\xspace}
\providecommand{\mbarn}{\ensuremath{\,\mathrm{mb}}\xspace}
\providecommand{\Mix}{\code{MixingModule}\xspace}
% written by minbias_catalogue.py: three tables of the minimum-bias library (do not edit)
\newcommand{\minbiasCatalogueTen}{%
\begin{tabular}{rlrrrrrrrrrrr}\toprule
event & proc. & interactions & stable & charged & central & $|\eta|>5.5$ & lead \pt & $\Sigma\pt$ & jets $>10$ & fwd $p$ & $z$ [cm] & $t$ [ns] \\ \midrule
1142 & ND & 1 & 74 & 32 & 2 & 27 & 0.88 & 21.0 & 0 & 1 & $-0.96$ & $-0.227$ \\
1143 & ND & 22 & 515 & 232 & 100 & 43 & 3.77 & 252.1 & 0 & 0 & $+1.10$ & $-0.006$ \\
1144 & ND & 23 & 511 & 240 & 62 & 86 & 4.17 & 225.4 & 1 & 0 & $-4.84$ & $+0.035$ \\
1145 & ND & 1 & 76 & 38 & 2 & 31 & 1.71 & 25.3 & 0 & 0 & $+0.90$ & $+0.150$ \\
1146 & ND & 28 & 560 & 264 & 102 & 63 & 4.11 & 288.5 & 0 & 0 & $+1.08$ & $+0.166$ \\
1147 & ND & 17 & 438 & 194 & 57 & 68 & 3.98 & 195.4 & 0 & 0 & $-1.60$ & $+0.096$ \\
1148 & SD & 1 & 32 & 12 & 0 & 19 & 1.85 & 11.3 & 0 & 1 & $-1.15$ & $-0.083$ \\
1149 & ND & 11 & 463 & 186 & 62 & 61 & 4.80 & 213.0 & 1 & 0 & $+0.97$ & $-0.044$ \\
1150 & ND & 19 & 345 & 148 & 37 & 59 & 4.49 & 154.7 & 0 & 0 & $-2.60$ & $+0.145$ \\
1151 & ND & 2 & 82 & 40 & 5 & 18 & 1.44 & 29.7 & 0 & 0 & $-2.61$ & $-0.073$ \\
\bottomrule\end{tabular}}
\newcommand{\minbiasCatalogueSummary}{%
\begin{tabular}{lrrrr}\toprule
per collision & mean & median & min & max \\ \midrule
interactions & 6.0 & 2 & 1 & 45 \\
stable particles & 179 & 116 & 3 & 829 \\
charged stable particles & 83 & 54 & 2 & 374 \\
central ($|\eta|<2.5$, $\pt>0.5\GeV$) & 19 & 8 & 0 & 167 \\
beyond $|\eta|=5.5$ & 37 & 34 & 2 & 108 \\
\pt of the hardest particle [GeV] & 2.17 & 1.7 & 0.1 & 17.9 \\
scalar sum of \pt [GeV] & 73 & 38 & 1 & 518 \\
\antikt $R=0.4$ jets above $10\GeV$ & 0.15 & 0 & 0 & 7 \\
forward protons ($\xi<0.55$) & 0.38 & 0 & 0 & 2 \\
vertex $z$ [cm], the 960 drawn & $0.22$ & $0.30$ & $-12.4$ & $13.9$ \\
vertex $t$ [ns], the 960 drawn & $0.003$ & $-0.001$ & $-0.35$ & $0.38$ \\
\bottomrule\end{tabular}}
\newcommand{\minbiasCataloguePerProcess}{%
\begin{tabular}{lrrrrrrrr}\toprule
per process & events & $\sigma$ [mb] & interactions & stable & charged & central & lead \pt [GeV] & $\Sigma\pt$ [GeV] \\ \midrule
non-diffractive (ND) & 1425 & 55.42 & 8.1 & 230 & 106 & 25 & 2.6 & 95 \\
single diffractive (SD) & 319 & 12.86 & 1.0 & 51 & 24 & 4 & 1.2 & 17 \\
double diffractive (DD) & 219 & 8.85 & 1.0 & 62 & 29 & 4 & 1.2 & 21 \\
central diffractive (CD) & 37 & 1.29 & 1.0 & 28 & 14 & 2 & 0.9 & 9 \\
\bottomrule\end{tabular}}

\chapter{The Crowd on the Road: Pileup}
\chaptermark{The Crowd on the Road}
\label{ch:crowd_on_the_road}

\begin{journeybox}
Our family is not alone. At every bunch crossing some sixty other families
set out at the same moment, and the road is so crowded that some of their
members walk in step with ours and are impossible to tell apart. The road
is the road of Run~3 in its last full year, 2024: the crowd of this
chapter is the one CMS recorded, and the one its simulation was built to
reproduce, and the family will meet it again at every stop ahead.
\end{journeybox}

\physicsline{the mean number of collisions per bunch crossing,
$\mup=L\sigmainel/(n_b f_{\mathrm{rev}})$, and its Poisson fluctuation; the
minimum-bias collision and its four kinds; in-time and out-of-time pileup,
the 25\ns clock and the sixteen crossings of the production window; the
beam spot of 2024, where and when the collisions happen; how \CMSSW mixes
a library of soft collisions into a signal event, why the production
premixes, and how a flat pileup profile is reweighted; the median energy
density \rhoPU, the jet area $A$ and the subtraction $\pt-\rhoPU A$. The
one number to take away: at $\mup=60$ the crowd puts $59\GeV$ into every
unit of area of the $(\eta,\phi)$ plane, thirty GeV of strangers into our
hundred-GeV jet, and the subtraction removes them on average, not event by
event: $\pm7\GeV$ stay.}

\section{What we will see at this stop}
\label{sec:crowd_on_the_road:what_we_compute}

At the end of \cref{ch:family_photograph} the photograph was taken: jet [0]
of \akfour, $99.83\GeV$ in $31$ particles around $\eta=-1.44$, $\phi=-2.71$,
the reference every later jet will be judged against. The family was alone
in the picture. It never is in the detector. The LHC does not collide one
proton with one proton; it collides bunches of $10^{11}$ protons every
$25\ns$, and at the luminosity the LHC delivered in 2024, about sixty pairs
of protons collided inelastically in each crossing. One of those collisions
is ours, the $q\bar q\to b\bar b$ of \cref{ch:ancestral_home}. The other
sixty happened because they could, soft and uneven, and their particles
land in the same detector at the same moment, on top of our family, inside
the circle of our jet. They are called pileup, and they are the crowd on
the road. This stop is about the crowd itself: how large it is and why,
what one member of it carries, where and when along the beam line it
walks, and how much of it falls inside our jet before any detector has
seen a single particle. The mitigation, the tracks that point back to the
right vertex, is the business of
\crefrange{ch:crossing_the_border}{ch:the_police}; the correction that
pays for what cannot be removed is \cref{ch:paying_for_uninvited}. Here we
only let the crowd arrive and measure what it does.

The crowd is the crowd of 2024, and the choice deserves a sentence. The
family of \crefrange{ch:ancestral_home}{ch:family_photograph} was born at
$\sqrt{s}=13.6\TeV$, the energy of Run~3, and Run~3 is over: its
luminosity, its filling scheme, its beam spot and its detector are
frozen, measured, and written into the conditions of \CMSSW, so that a
book pinned to them stays right. The High-Luminosity LHC, with its two
hundred collisions per crossing and the detector CMS is building for
them, is still being designed, and a book pinned to it would be out of
date before the first beam. So the crowd of every stop from here on is
the Run-3 crowd, the detector of \cref{ch:crossing_the_border} is the
2024 CMS detector, and what the HL-LHC will change is said at the end of
each chapter, in its open problems, where the scripts of this stop keep
the HL-LHC road one option away.\footnote{\file{standalone_pileup.py
--scenario ultimate} mixes the same family at $\mup=200$ with the Phase-2
luminous region, and the exercises and open problems use it.}

The numbers to take away are four. The crowd of our crossing has $63$
collisions, a Poisson draw with mean $60$, and brings $12\,113$ particles
to the $744$ of the family, $860$ charged ones into the tracker's reach
against $93$. It puts $68\GeV$ into every unit of area of the
$(\eta,\phi)$ plane in this crossing, $59$ on average, $0.94$ per
collision, so that the circle of our jet, $0.53$ units of area, receives
$33.61\GeV$ of strangers from $34$ different collisions and reads
$132.75\GeV$ instead of $99.83$. The subtraction $\pt-\rhoPU A$ takes
$36.09\GeV$ back out and leaves $96.66$, three GeV below the truth in this
crossing and $-2\pm7\GeV$ from it over four hundred draws of the crowd:
the average is removed, the fluctuation is not, and that fluctuation is
the first term of the resolution of a jet in Run~3. And the collisions are
spread over $\pm3.6\cm$ along the beam and $\pm120\ps$ in time, the
nearest one to ours a twentieth of a millimetre away in this crossing and
a millimetre on the median, which is the room the vertex finder will have
to work in, with no timing layer to help it. As at the previous stops we
meet the numbers four times, from the top down: a film shows the
photograph of \cref{ch:family_photograph} disappearing under the crowd
collision by collision, the beam spot and the 25\ns clock, one member of
the crowd, the median that measures it, our jet filling up and being
subtracted, and the crowd drawn four hundred times; a Python file, needing
only NumPy, mixes a library of two thousand minimum-bias collisions into
the record in the way \CMSSW's \Mix does, clusters the result with the
\antikt of \cref{ch:family_photograph} and ghosts, and writes \rhoPU, the
areas and the subtracted jets; the theory derives \mup from the
luminosity, the collision from the cross section, the beam spot from two
bunches and \rhoPU from the median; and the official packages, \PYTHIA for
the library and its cross sections, \FASTJET for the areas and the two
estimators of \rhoPU on exactly our mixed particles, and the 2024 \CMSSW
mixing of the companion, get the last word.

\paragraph{The film.}
\label{sec:crowd_on_the_road:animation}
The film first: the eighth on the playlist of \jitjetworld puts the
photograph into the crowd in six narrated scenes, every number on the
screen read from the record and the ensemble that the script wrote. The
voice picks up where the film of \cref{ch:family_photograph} stopped, the
photograph is taken, and ends at the next stop, the border of the
detector. The $651$ signal particles sit in the $(\eta,\phi)$ plane with
their jet circles, and the $63$ collisions of the crossing stack onto
them one by one, a counter running and \rhoPU rising by a GeV per
collision until the plane is khaki and \antikt finds $173$ jets; the beam
spot appears as the overlap of two bunches, the vertices of our crossing
spread along it with the signal among them, and the fifteen neighbouring
crossings of the production window arrive $25\ns$ apart; the four kinds
of minimum-bias collision are drawn one after the other, each with its
row in the table of cross sections and multiplicities, and the
multiplicity of all two thousand builds up; the grid cells fill
with transverse momentum and their median settles at \rhoPU; the family's
circle fills with strangers, the counter climbs from $99.8$ to
$132.7\GeV$, and $\rhoPU A$ is taken out; and the crowd is drawn four
hundred times on the same photograph, the pileup inside the circle and
the residual after the subtraction filling their histograms.
\Cref{fig:crowd_on_the_road:stills} shows one key frame per scene. Watch
first, and let the formulae catch up.

\begin{figure}[htbp]\centering
\includegraphics[page=1,width=0.325\textwidth]{PU/Standalone/animate_standalone_pileup_stills.pdf}\hfill
\includegraphics[page=2,width=0.325\textwidth]{PU/Standalone/animate_standalone_pileup_stills.pdf}\hfill
\includegraphics[page=3,width=0.325\textwidth]{PU/Standalone/animate_standalone_pileup_stills.pdf}\\[1.5mm]
\includegraphics[page=4,width=0.325\textwidth]{PU/Standalone/animate_standalone_pileup_stills.pdf}\hfill
\includegraphics[page=5,width=0.325\textwidth]{PU/Standalone/animate_standalone_pileup_stills.pdf}\hfill
\includegraphics[page=6,width=0.325\textwidth]{PU/Standalone/animate_standalone_pileup_stills.pdf}
\caption[Key frames of the animated standalone pileup stage]{Key frames of
the six scenes of the animated standalone pileup stage, drawn by
\file{animate_standalone_pileup.py} with \code{--stills}, in reading
order. The crowd arriving on the photograph; the beam spot and the
sixteen crossings on the clock; the four kinds of collision and the library;
the grid cells and their median; our jet filling up and being
subtracted; the crowd drawn four hundred times. The film itself is the
eighth on the playlist of \jitjetworld.}
\label{fig:crowd_on_the_road:stills}
\end{figure}

\paragraph{The picture to carry: the crowd.}
The two halves of \cref{fig:crowd_on_the_road:crowd_sketch} are the whole
chapter. On the left, one bunch crossing: two bunches of protons pass
through each other, and wherever and whenever they overlap a collision
can happen, sixty of them in a region $4\cm$ long and an eighth of a
nanosecond wide, one of them ours. Nothing selects it; the detector
records the crossing because our collision fired the trigger, and the
other sixty come free, as do the crowds of the twelve crossings before
and the three after, which the slow detectors remember. On the right, the
circle of our jet in the $(\eta,\phi)$ plane: the family, thirty-one
particles within a tenth of a unit of the axis, and the strangers,
sixty-five soft particles from thirty-four collisions scattered evenly
over the circle, a third as many GeV as the family itself. Two things in
the picture matter for everything that follows. The crowd is
\emph{uniform on average and uneven in every event}: its mean density
\rhoPU is one number for the whole plane and can be measured from the
empty parts of the plane and subtracted from the full ones, and that
subtraction is exact only for the mean, because the number of collisions,
the number of particles each brings and where they land all fluctuate,
and what they leave behind in a circle of area $0.5$ is $\pm7\GeV$ wide.
And the crowd is \emph{soft}: its hardest particle is $2\GeV$ on average,
its members make jets of their own above $30\GeV$ only before \rhoPU is
subtracted and never after, and what they do to our jet is a pedestal,
not a disguise; a fake jet of the crowd is rare, a true jet with the
crowd on top is every jet. The rules are the jet-area paper of Cacciari
and Salam~\cite{Cacciari:2007fd} and the catchment area of
jets~\cite{Cacciari:2008gn}, the \FASTJET manual~\cite{Cacciari:2011ma},
the \CMSSW \Mix and its \code{BetafuncEvtVtxGenerator}; the experiment's
view of the crowd is the pileup mitigation paper of
CMS~\cite{CMS:2020ebo}, and the theorist's the review of
Soyez~\cite{Soyez:2018opl}.

\begin{figure}[htbp]\centering
\resizebox{\textwidth}{!}{% ---------------------------------------------------------------------------
% The picture to carry: (a) one bunch crossing of Run 3 (2025), sixty
% collisions along the beam spot, ours among them; (b) the circle of our jet
% in the (eta, phi) plane: the family and the strangers that walked in, with
% the numbers of the 400-draw ensemble of standalone_pileup.py --repeat.
% ---------------------------------------------------------------------------
\begin{tikzpicture}[x=1cm,y=1cm,
  every node/.style={font=\footnotesize},
  lab/.style={font=\scriptsize,align=center},
  head/.style={font=\small\bfseries},
  fam/.style={circle,fill=red!80!black,inner sep=0pt,minimum size=4pt},
  crowd/.style={circle,fill=olive!60!brown,inner sep=0pt,minimum size=2.6pt},
  vtx/.style={circle,fill=cyan!60!blue,inner sep=0pt,minimum size=2.6pt},
  ring/.style={red!80!black,line width=1.2pt}]
% ---------------- (a) the bunch crossing ----------------------------------------------------------
\begin{scope}
\node[head,anchor=west] at (-0.3,3.2) {(a) one bunch crossing of 2025: $\mu=60$ collisions along $z$};
% the two bunches moving towards each other
\fill[blue!25] (-4.2,0.9) ellipse (1.6 and 0.32);
\fill[orange!35] (4.2,0.9) ellipse (1.6 and 0.32);
\draw[-{Stealth[length=2mm]},blue!60!black,thick] (-2.4,1.35) -- (-1.4,1.35) node[lab,above,pos=0.5] {beam 1};
\draw[-{Stealth[length=2mm]},orange!80!black,thick] (2.4,1.35) -- (1.4,1.35) node[lab,above,pos=0.5] {beam 2};
\node[lab,text=blue!60!black] at (-4.2,0.45) {$10^{11}$ protons per bunch};
\node[lab,text=orange!80!black] at (4.2,0.45) {every 25 ns};
% the beam spot: a Gaussian along z, sigma_z = 3.6 cm (scale: 0.25 cm per cm)
\draw[-{Stealth}] (-4.5,-0.3) -- (4.5,-0.3) node[below left,lab] {$z$ [cm]};
\foreach \x/\t in {-4/$-16$,-2/$-8$,0/0,2/8,4/16} {\draw (\x,-0.35) -- (\x,-0.25); \node[lab,anchor=north] at (\x,-0.4) {\t};}
\draw[green!50!black,thick] plot[domain=-4.5:4.5,samples=80] (\x,{-0.3+1.1*exp(-(\x-0.1)*(\x-0.1)/(2*0.9*0.9))});
\node[lab,text=green!50!black,anchor=west] at (1.2,0.75) {$\sigma_z=3.6$ cm, $\sigma_t=120$ ps};
% 60 vertices drawn along z (pseudo-random, Gaussian-like by summing uniforms)
\pgfmathsetseed{8}
\foreach \i in {1,...,60} {
  \pgfmathsetmacro{\zz}{0.1+0.9*(rand+rand+rand)*0.58}
  \pgfmathsetmacro{\yy}{-0.3+0.08+0.5*rnd*exp(-\zz*\zz/(2*0.9*0.9))}
  \node[vtx] at (\zz,\yy) {};
}
\node[star,star points=5,fill=red!80!black,inner sep=0pt,minimum size=9pt] at (-0.05,0.05) {};
\node[lab,text=red!80!black,anchor=north] at (-0.05,-0.05) {ours};
\node[lab,anchor=north west,text width=9.2cm,align=left] at (-4.5,-0.95) {Every 25 ns two bunches cross; $\mu=L\sigma_{\rm inel}/(n_b f_{\rm rev})\simeq60$ inelastic collisions happen in the same crossing, spread over $\pm10$ cm and $\pm0.35$ ns. One is ours. The others are the crowd: soft, uneven, and in the same picture. The twelve crossings before and the three after bring their own crowds, $25$ ns apart.};
\end{scope}
% ---------------- (b) the circle of our jet -------------------------------------------------------
\begin{scope}[xshift=11cm]
\node[head,anchor=west] at (-0.5,3.2) {(b) the circle of our jet, $R=0.4$, with the crowd inside};
\draw[ring] (1.6,0.9) circle (1.7);
\node[lab,text=red!80!black,anchor=south west] at (2.75,2.15) {$R=0.4$};
% the family: 31 particles, collimated
\pgfmathsetseed{21}
\foreach \i in {1,...,31} {
  \pgfmathsetmacro{\r}{1.7*(0.08+0.55*rnd*rnd)}
  \pgfmathsetmacro{\a}{360*rnd}
  \node[fam] at ({1.6+\r*cos(\a)},{0.9+\r*sin(\a)}) {};
}
% the crowd: about 130 soft particles per unit of area, 65 of them inside the circle
\pgfmathsetseed{34}
\foreach \i in {1,...,165} {
  \pgfmathsetmacro{\xx}{-0.9+5.0*rnd}
  \pgfmathsetmacro{\yy}{-1.4+4.6*rnd}
  \node[crowd] at (\xx,\yy) {};
}
\node[lab,anchor=north west,text width=5.9cm,align=left] at (4.1,3.0) {\textcolor{red!80!black}{\textbf{the family}}: 31 particles, $99.8$ GeV, $70$ of them inside $0.1$ of the axis\\[3pt]
  \textcolor{olive!60!brown}{\textbf{the crowd}}: $\rho\simeq59$ GeV per unit area at $\mu=60$, $0.97$ per collision; the circle covers $\pi R^2=0.50$, so $28\pm8$ GeV walk in, 68 particles from 30 collisions\\[3pt]
  \textbf{the subtraction}: $p_T-\rho A$ takes $29\pm5$ out; what is left fluctuates by $\pm6$ GeV around the family's $100$\\[3pt]
  \textbf{the one number}: thirty GeV of strangers on a hundred-GeV jet, removed on average, not event by event};
\end{scope}
\end{tikzpicture}
}
\caption[The crowd on the road: one bunch crossing, and the circle of our
jet with the strangers inside]{The crowd on the road. Left: one bunch
crossing of 2024, two bunches passing through each other and the sixty
collisions of the crossing spread along their overlap, ours among them.
Right: the circle of our jet, the family in red and the crowd in khaki,
with the numbers of the four hundred draws of the crowd on the same
photograph: what walks in, what the subtraction takes out, and what is
left.}
\label{fig:crowd_on_the_road:crowd_sketch}
\end{figure}

\paragraph{The number, four ways.}
\label{sec:crowd_on_the_road:validation}
\label{sec:crowd_on_the_road:standalone_code}
\Cref{tab:crowd_on_the_road:numbers} gives away the ending, and its columns
are of two kinds. The crowd itself is made by \PYTHIA: a library of two
thousand minimum-bias collisions at $13.6\TeV$, generated once with the
fragment the Run-3 libraries use, \code{SoftQCD:inelastic} with the CP5
tune, and stored as a record, so that the standalone script needs NumPy
and nothing else to read it. That is the way of \CMSSW too, whose \Mix
reads a library of simulated collisions made once and stacked many
times, and the first rows of the table check the library against the
cross sections \PYTHIA printed and against the $250$-event library of the
2024 \CMSSW exercise of the companion. Everything after the library is a
port: the Poisson draw of the \Mix, the beam spot of the
\code{BetafuncEvtVtxGenerator}, the stacking, the \antikt with ghosts, the
two estimators of \rhoPU and the subtraction; and for those the sharpest
check is, as in \cref{ch:family_photograph}, a comparison on identical
input, the $9684$ particles of the mixed final state clustered by
\FASTJET~3.5.2 with its active area and its two background estimators.
The $173$ jets agree exactly, to the last constituent and the last digit
of every momentum; the grid median agrees to the last digit, since the
grid is the same grid; the areas agree to the placement of the ghosts,
which \FASTJET scatters differently from the script, $\pm0.03$ on an area
of $0.5$; and the jet median, which is measured on the ghost jets
themselves, agrees to the percent by which it moves when the ghosts move.
The last column is the 2024 \CMSSW mixing of the companion,
\code{CMSSW\_15\_0\_5} with the 2024 conditions at $13.6\TeV$: the GEN-SIM
library it reads, its vertex smearing, and one signal event mixed at
$\mup=60$ in three crossings, which is the exercise a laptop can afford
once the hits of \GEANT are involved, and which the standalone script
reproduces at the generator level with \code{--window 1 --sequential}.

The standalone script writes out what the \Mix does before the
digitisers run and what the \FASTJET area tools do afterwards. Its
learning mode, \code{-v}, prints the arithmetic of \mup and the sixteen
Poisson draws, which library events were read, the payload of the beam
spot and the vertices drawn, the stacking crossing by crossing, the
clustering of our jet with its signal and pileup constituents, the two
estimates of \rhoPU, the subtraction of the twelve leading jets against
the truth, the vertices the tracker can see and the profile weights;
\code{-vv} adds every collision with its vertex and every constituent of
our jet. The \code{-v} log for seed~1 is committed, and every number
quoted in the theory of this chapter is taken from it. How each column of
the table was produced, the commands, the modes and the logs, is in the
companion of this chapter,
\begin{center}
  \companion{chapter_08_PU}
\end{center}
together with the full tables behind every number quoted below.

\begin{table}[htbp]\centering\small
\caption[The crowd on the road obtained four ways]{The crowd on the road
obtained four ways. ``Library'' is the two thousand minimum-bias
collisions of \file{pileup_library.lhe.gz}, made once with \PYTHIA at
$13.6\TeV$; ``standalone'' is \file{standalone_pileup.py} on the record of
\cref{ch:family_photograph} with that library, seed~1, $\mup=60$ in
$\bx=-12\ldots+3$; ``\FASTJET'' is \FASTJET~3.5.2 on exactly the mixed
final state of the record; ``400 draws'' is the in-time crowd drawn again
and again on the same photograph, with the pileup counted inside the
circle of $R=0.4$ around our jet; ``\CMSSW'' is the 2024 library and
mixing of the companion (\code{CMSSW\_15\_0\_5}, $13.6\TeV$, $250$ library
events, one signal event mixed at $\mup=60$ in $\bx=-1\ldots+1$). The
spread of a single event follows $\pm$.}
\label{tab:crowd_on_the_road:numbers}
\scriptsize\setlength{\tabcolsep}{3pt}
\resizebox{\textwidth}{!}{%
\begin{tabular}{lrrrrr}\toprule
 & \PYTHIA & library & standalone / \FASTJET & 400 draws & \CMSSW 2024 \\
 & cross sections & 2000 collisions & seed 1, $\mup=60$ & $\mup=60$ & 250 + 1 events \\ \midrule
\sigmainel at $13.6\TeV$ [mb] & $78.42$ & -- & -- & -- & $78.42$ \\
ND / SD / DD / CD [\%] & $70.7$ / $16.4$ / $11.3$ / $1.6$ & $71.2$ / $16.0$ / $11.0$ / $1.9$ & $69.5$ / $16.9$ / $12.0$ / $1.7$ (960 drawn) & -- & $73$ / $18$ / $8$ / $1$ \\
particles per collision: mean, median & -- & $179$, $116$ & $192$ & -- & $194$, $132$ \\
\quad charged, $|\eta|<2.5$, $\pt>0.5\GeV$ & -- & $12.2$ (median 6) & $13.7$ & -- & $21$ (all, median 10) \\
\quad hardest particle [GeV] & -- & $2.2$ (max $17.9$) & -- & -- & $2.2$ (max $8.1$) \\ \midrule
\mup: formula / mixed & -- & -- & $59.3$ / $60$ & $60$ & $60$ (AVE\_60); production flat $55$--$75$ \\
collisions per crossing & -- & -- & $50$ to $68$ in 16 crossings, $\mathbf{63}$ in time, $960$ & $59.9\pm7.4$ (in time) & $66$, $\mathbf{55}$, $77$ ($\bx=-1\ldots+1$) \\
beam spot $\sigma_z$ [cm], $\sigma_t$ [ps] & -- & -- & payload $3.60$, $120$; drawn $3.79$, $119$ & $3.74$, $119$ & $3.6$, $120$ (payload); $3.56$, $122$ (250) \\
signal vertex $z$ [cm], $t$ [ps] & -- & -- & $-0.17$, $-121$ & -- & $3.11$, $-140$ \\
nearest pileup vertex: $\Delta z$ [mm], $\Delta t$ [ps] & -- & -- & $0.06$, $5$ & median $0.98$, $2.6$ & $0.4$ (3 within 1 mm) \\ \midrule
in-time particles, within $|\eta|<5$ no $\nu$ & -- & -- & $12\,113$, $9033$ & $7949$ & $31\,597$ primaries (198 ev.) \\
\rhoPU, grid median [GeV per unit area] & -- & -- & $68.28$ / 68.28 & $59.0\pm11.5$ & -- \\
\rhoPU, $k_T$-jet median & -- & -- & $69.87$ / 70.43 & -- & -- \\
\quad per collision & -- & $0.95$ (plateau) & $1.08$ / $1.11$ & $0.94$ (slope) & -- \\
mixed jets above $3$ / $30\GeV$ & -- & -- & $173$ / $62$, all identical & -- & -- \\
our jet: \pt with the crowd [GeV], $n$ & -- & -- & $132.75$, $112$ (31 + 81) & -- & -- \\
\quad area $A$ & -- & -- & $0.529$ / 0.529 & $\pi R^2=0.503$ & -- \\
\quad pileup inside, $\rhoPU A$, $\pt-\rhoPU A$ [GeV] & -- & -- & $33.61$, $36.09$, $96.66$ & $27.8\pm8.4$, $29.7\pm5.8$ & $\sim55$ (ECAL energy, in time) \\
\quad residual to the signal inside [GeV] & -- & -- & $-3.41$ & $-1.9\pm6.7$ & -- \\
pileup jets above $30\GeV$: before, after & -- & -- & $56$, $0$ & -- & -- \\
in-time collisions visible to the tracker & -- & -- & $43$ of $63$ & $62\%$ & -- \\
\bottomrule
\end{tabular}}
\end{table}

\section{Let us dive into the details: the theory}
\label{sec:crowd_on_the_road:theory}

\subsection{How many: the luminosity, the cross section and the Poisson draw}
\label{sec:crowd_on_the_road:forty_families_at_once_what_pileup_is_an}

The question first. The hard scattering of \cref{ch:ancestral_home} has a
cross section of $345\pb$; an inelastic proton-proton collision of any
kind has one of $78\mbarn$, eleven orders of magnitude more. Every time
two bunches cross, each proton of the one sees the $10^{11}$ protons of
the other, and the number of inelastic collisions in that crossing is
the luminosity of the crossing times \sigmainel. Nothing about our
collision makes the others less likely; they are there whenever the
beams are, and the only question is how many.
\Cref{fig:crowd_on_the_road:mu_sketch} draws the answer. The luminosity
$L$ an accelerator quotes is a rate per unit cross section integrated
over all crossings, so the number of inelastic collisions per second is
$L\sigmainel$; the number of bunch crossings per second is the number of
colliding bunch pairs times the revolution frequency, and the ratio is
the mean number of collisions per crossing:
\begin{equation}
  \mup=\frac{L\,\sigmainel}{n_b\,f_{\mathrm{rev}}},\qquad
  f_{\mathrm{rev}}=\frac{c}{26\,659\,\mathrm{m}}=11\,245\,\mathrm{Hz},\qquad
  n_b=2350\ \text{of the }3564\text{ slots in }2024.
  \label{eq:crowd_on_the_road:mu}
\end{equation}
The denominator is not the $40\,\mathrm{MHz}$ of the $25\ns$ spacing:
only $2350$ of the $3564$ slots of the ring held a colliding pair in the
filling scheme of 2024, the gaps being needed by the injection and dump
kickers, so that CMS saw $2.64\times10^{7}$ crossings per second. The
numbers of 2024 are $L=2.0\times10^{34}\,\mathrm{cm^{-2}s^{-1}}$ at the
levelled peak and, from \PYTHIA with the CP5 tune at $13.6\TeV$,
$\sigmainel=78.42\mbarn$: the script prints $1.57\times10^{9}$
collisions per second over $2.64\times10^{7}$ crossings per second,
$\mup=59.3$, or $60.5$ with the $80\mbarn$ that CMS uses for its pileup
bookkeeping. The mixing scenario of the 2024 simulation rounds this to
$60$, \code{AVE\_60\_BX\_25ns}, and $60$ is the number this chapter mixes;
the same arithmetic with the $2748$ bunches of the HL-LHC gives $127$ for
its nominal $5\times10^{34}$ and $190$ for its ultimate $7.5\times10^{34}$,
which \CMSSW mixes as $140$ and $200$.

\begin{figure}[htbp]\centering
\resizebox{\textwidth}{!}{% ---------------------------------------------------------------------------
% How many: mu = L sigma_inel / (n_b f_rev). (a) the ring: 3564 slots of 25 ns,
% 2340 of them colliding at CMS in 2025; (b) the arithmetic of the 2025
% scenario (the header of the CMSSW 2025 profile) against the HL-LHC ones; (c) the Poisson draw and the flat profile
% of the Run-3 production.
% ---------------------------------------------------------------------------
\begin{tikzpicture}[x=1cm,y=1cm,
  every node/.style={font=\footnotesize},
  lab/.style={font=\scriptsize,align=center},
  head/.style={font=\small\bfseries},
  bunch/.style={circle,fill=blue!70!black,inner sep=0pt,minimum size=2.4pt},
  gap/.style={circle,draw=black!30,inner sep=0pt,minimum size=2.4pt}]
% ---------------- (a) the ring ---------------------------------------------------------------------
\begin{scope}
\node[head,anchor=west] at (-2.4,3.0) {(a) the ring: bunches every 25 ns};
\draw[black!50,thick] (0,0) circle (2.1);
\foreach \i in {0,...,71} {
  \pgfmathsetmacro{\a}{5*\i}
  \pgfmathparse{int(mod(\i,9))}
  \ifnum\pgfmathresult<6 \node[bunch] at ({2.1*cos(\a)},{2.1*sin(\a)}) {};
  \else \node[gap] at ({2.1*cos(\a)},{2.1*sin(\a)}) {};
  \fi
}
\node[star,star points=5,fill=red!80!black,inner sep=0pt,minimum size=9pt] at (2.1,0) {};
\node[lab,anchor=west] at (2.3,0) {CMS};
\node[lab,text=blue!70!black] at (0,0.5) {26\,659 m, $f_{\rm rev}=11\,245$ Hz};
\node[lab] at (0,0.0) {3564 slots of 25 ns,};
\node[lab] at (0,-0.4) {2340 pairs collide at CMS\\ (the gaps are for the kickers)};
\node[lab,anchor=north,text width=4.6cm] at (0,-2.3) {bunch crossings per second at CMS in 2025: $n_b f_{\rm rev}=2340\times11\,245=2.63\times10^{7}$, not the 40 MHz of the spacing};
\end{scope}
% ---------------- (b) the arithmetic ----------------------------------------------------------------
\begin{scope}[xshift=6.6cm]
\node[head,anchor=west] at (-0.4,3.0) {(b) collisions per second over crossings per second};
\node[anchor=north west,text width=9.2cm,align=left,font=\footnotesize] at (-0.4,2.5) {%
  $\displaystyle \mu=\frac{L\,\sigma_{\rm inel}}{n_b\,f_{\rm rev}}
   =\frac{2.23\times10^{34}\,{\rm cm^{-2}s^{-1}}\times 78.4\times10^{-27}\,{\rm cm^2}}{2340\times 11\,245\,{\rm s^{-1}}}$\\[3pt]
  \hspace*{9mm}$\displaystyle =\frac{1.75\times10^{9}\ {\rm collisions/s}}{2.63\times10^{7}\ {\rm crossings/s}}=66.5$\\[6pt]
  {\scriptsize\begin{tabular}{@{}lcccc@{}}
   scenario & $L$ [cm$^{-2}$s$^{-1}$] & $n_b$ & $\mu$ (formula) & mixed at \\ \hline
   Run 3, 2025, levelled & $2.23\times10^{34}$ & 2340 & 66 & \textbf{60} \\
   HL-LHC nominal & $5\times10^{34}$ & 2748 & 127 & 140 \\
   HL-LHC ultimate & $7.5\times10^{34}$ & 2748 & 190 & 200 \\
  \end{tabular}}\\[4pt]
  $L$ and $n_b$ of 2025 are the header of the \CMSSW 2025 profile (levelled for 9 of 12.5 hours; mean 64 over the year). The book mixes at \code{AVE\_60\_BX\_25ns}, as its \CMSSW companion; the production draws $\mu_{\rm true}$ flat in $55\ldots75$ and reweights to the measured profile. $\sigma_{\rm inel}$ is \PYTHIA's with CP5 ($\sigma_{\rm tot}=100.3$, $\sigma_{\rm el}=21.9$ mb); the 80 mb of the bookkeeping gives 67.8.};
\end{scope}
% ---------------- (c) the Poisson draw --------------------------------------------------------------
\begin{scope}[xshift=16.6cm]
\node[head,anchor=west] at (-0.4,3.0) {(c) the draw: $n\sim{\rm Poisson}(\mu)$, $\mu\pm\sqrt{\mu}$};
\draw[-{Stealth}] (0,0) -- (5.2,0) node[below left,lab] {$n$};
\draw[-{Stealth}] (0,0) -- (0,2.4) node[left,lab] {$P(n)$};
% x = 0.3 + (n - 30) * 0.08: n = 30 at 0.3, n = 90 at 5.1
\foreach \x/\t in {1.1/40,2.7/60,4.3/80} {\draw (\x,-0.05) -- (\x,0.05); \node[lab,anchor=north] at (\x,-0.08) {\t};}
% the production's flat profile of mu_true, 55 ... 75
\fill[black!12] (2.3,0) rectangle (3.9,0.55);
\node[lab,text=black!60,anchor=west] at (4.0,0.85) {the production:\\ $\mu_{\rm true}$ flat in $55\ldots75$};
% Poisson(60) ~ Gaussian of width 7.7: 0.62 cm, height 2.0 at n = 60
\draw[green!50!black,thick] plot[domain=0.3:5.1,samples=70] (\x,{2.0*exp(-(\x-2.7)*(\x-2.7)/(2*0.62*0.62))});
\node[lab,text=green!50!black,anchor=west] at (3.4,2.1) {$\mu=60$: $60\pm8$ (2025)};
\node[lab,anchor=north west,text width=6.2cm,align=left] at (-0.4,-0.75) {For every crossing of the window $bx=-12\ldots+3$ the mixing draws its own $n$: sixteen numbers around 60, 960 collisions per signal event. The \code{PileupSummaryInfo} keeps them with the mean (\code{getTrueNumInteractions}), which the reweighting reads.};
\end{scope}
\end{tikzpicture}
}
\caption[How many: the ring, the arithmetic of $\mup$ and the Poisson
draw]{How many. (a) The ring: $3564$ slots of $25\ns$, $2350$ of them
colliding at CMS in 2024, $11\,245$ times a second each. (b) The
arithmetic of \cref{eq:crowd_on_the_road:mu} for 2024 and for the two
HL-LHC scenarios, and the value \CMSSW mixes at. (c) The draw: for every
crossing of the window the \Mix takes $n$ from a Poisson distribution of
mean \mup, sixteen numbers around $60$ for our event, and the flat band
of \mutrue the production draws from.}
\label{fig:crowd_on_the_road:mu_sketch}
\end{figure}

\mup is a mean. The collisions of a crossing are independent, so their
number is a Poisson variable with mean \mup and standard deviation
$\sqrt{\mup}$, $60\pm8$, and this is exactly what the \Mix draws, once
for every bunch crossing of the window it is asked to fill, $\bx=-12$ to
$+3$ in the Run-3 production, sixteen draws and $960$ collisions per
signal event. The window is asymmetric because the detectors with the
longest memory, the calorimeter pulses and the strip readout, integrate
collisions up to $300\ns$ old, while nothing arrives from the future
beyond the $75\ns$ of the latest pulse tails. The script draws the same
way, and seed~1 gives sixteen numbers from $50$ to $68$, $960$ collisions
for an expectation of $960$; the one that matters most, the in-time
crossing $\bx=0$, has $63$, and the draws are written to the record the
way \CMSSW writes them to its \code{PileupSummaryInfo}, one entry per
crossing with the number drawn and the mean it was drawn from. For every
collision the script then reads one event of the library at a random
position, with replacement, which is \CMSSW's default
(\code{sequential = False}) and the reason a library of two thousand can
serve a crossing of sixty and an ensemble of four hundred crossings:
$763$ distinct events served the $960$ collisions of our event, and
reusing a soft collision as the background of two unrelated signal events
costs nothing a measurement can see, as long as it is not reused within
one event. \Cref{fig:crowd_on_the_road:theory_mu} draws
\cref{eq:crowd_on_the_road:mu} against the luminosity with the three
scenarios on it, the Poisson distributions of $\mup=60$ and $200$ with
the sixteen draws of our event, and the third thing a mean hides.

\subsubsection{The profile of a year, and the reweighting}
\label{sec:crowd_on_the_road:profile}
\mup is a mean over a crossing, but it is also a number that changes
during a fill and from fill to fill: a year of running is a distribution
of \mup, the \emph{pileup profile}, weighted by the collisions themselves
since every collision counts once. In 2024 the luminosity was
\emph{levelled}: the beams were kept slightly separated at the start of a
fill so that \mup stayed near $60$ for hours instead of starting higher
and decaying, and only the last hours of a fill ran below it. The script
builds the profile from a luminosity curve: one fill levelled at $60$ for
eight hours and then decaying with a lifetime of ten keeps $72\%$ of its
collisions at the levelled value and has a mean of $57$; a year of such
fills levelled between $57$ and $67$, a model of the fill-to-fill spread
and not the measured profile, has a mean of $59$ with two thirds of its
collisions between $52$ and $66$. A simulated sample is mixed with one
profile and compared with data taken with another, and the repair is a
weight per event,
\begin{equation}
  w(\mutrue)=\frac{P_{\mathrm{target}}(\mutrue)}{P_{\mathrm{sample}}(\mutrue)},
  \label{eq:crowd_on_the_road:reweight}
\end{equation}
with \mutrue the mean the event's in-time crossing was drawn from, the
\code{getTrueNumInteractions} of the \code{PileupSummaryInfo}, and never
the number actually drawn, which would fold the Poisson fluctuation into
the weight twice. This is why the Run-3 production does not mix at a
fixed $60$ at all: its scenario, \code{Run3\_Flat55To75\_PoissonOOTPU},
hands the \Mix a histogram in place of a mean (\code{type = probFunction}),
flat between $55$ and $75$, draws \mutrue from it for every signal event,
and gives the sixteen crossings Poisson draws of that mean, about $1040$
collisions per signal event. A flat profile covers whatever the year
turns out to be, and the reweighting to the measured profile then costs
events where the band is wider than the year: against the model year the
weights run from $1.6$ at the levelled values to zero above $70$, $1.54$
at $\mutrue=60$. A sample mixed at a fixed $\mup=60$, like the record of
this chapter, covers one value and cannot be reweighted to a profile at
all. Panel (c) of \cref{fig:crowd_on_the_road:theory_mu} shows the three
profiles and the weight.

\begin{figure}[htbp]\centering
\includegraphics[page=1,width=\textwidth]{PU/Standalone/theory_pileup.pdf}
\caption[The crowd's size: $\mup$ against the luminosity, the Poisson draw,
the profiles]{Page~1 of \file{theory_pileup.pdf}. (a)
\Cref{eq:crowd_on_the_road:mu} against the luminosity, with the three
scenarios from the formula and the values \CMSSW mixes at. (b) The Poisson
distributions of $\mup=60$ and $200$, and the sixteen draws of our event,
one per crossing. (c) The pileup profile of a 2024-like year of levelled
fills, one fill, the flat $55$--$75$ profile of the Run-3 production, and
the weight of \cref{eq:crowd_on_the_road:reweight} for an event of the
flat sample.}
\label{fig:crowd_on_the_road:theory_mu}
\end{figure}

\subsection{One member of the crowd: the minimum-bias collision}
\label{sec:crowd_on_the_road:one_member}

What does one of the sixty carry? Not a hard scattering: a collision with
a $100\GeV$ parton pair in it is one in a hundred thousand, and the crowd
is what happens in all the others, the inelastic collisions in their
natural proportions, which the experiments call \emph{minimum bias} after
the trigger that selects them with the least possible bias. \PYTHIA makes
them as \code{SoftQCD:inelastic}, the sum of the four processes of
\cref{fig:crowd_on_the_road:minbias_sketch}: the non-diffractive
collision, in which both protons break up and several parton pairs
interact as in \cref{ch:uninvited_cousins}, so that the plateau of soft
particles stretches over ten units of rapidity; the single-diffractive
one, in which one proton survives and the other dissociates, leaving a
gap in rapidity on the surviving side; the double-diffractive one, both
dissociating with a gap between them; and the central-diffractive one,
both surviving and a small system in the middle. Their cross sections are
not computed from the parton distributions: the CP5 tune fixes the total
and the elastic cross sections, $\sigma_{\mathrm{tot}}=100.31$ and
$\sigma_{\mathrm{el}}=21.89\mbarn$, at whatever energy it is run, the
diffractive pieces come from \PYTHIA's parametrisation and the
non-diffractive part is what remains,
\begin{equation}
  \sigmainel=\sigma_{\mathrm{tot}}-\sigma_{\mathrm{el}}=78.42\mbarn
  =\underbrace{55.42}_{\mathrm{ND}}+\underbrace{2\times6.43}_{\mathrm{SD}}+\underbrace{8.85}_{\mathrm{DD}}+\underbrace{1.29}_{\mathrm{CD}}\ \mbarn,
  \label{eq:crowd_on_the_road:sigma}
\end{equation}
the shares $70.7$, $16.4$, $11.3$ and $1.6\%$, which the library
reproduces as $71.2$, $16.0$, $11.0$ and $1.9$ in its two thousand
collisions and the $250$-event \CMSSW library as $183$, $45$, $19$ and
$3$. The $14\TeV$ fragment of the release gives the same numbers to the
second decimal, which is why the Run-3 workflows can mix their
$13.6\TeV$ signals with a $14\TeV$ library and why the companion sets the
energy by hand instead. The measured inelastic cross section at $13\TeV$
is $79.5\pm1.8\mbarn$~\cite{ATLAS:2016ygv,CMS:2018mlc}, and the $80\mbarn$
of the bookkeeping is that number rounded.

\begin{figure}[htbp]\centering
\resizebox{\textwidth}{!}{% ---------------------------------------------------------------------------
% One member of the crowd: the four inelastic process types of SoftQCD:inelastic
% with their rapidity profiles (ND: both protons break up, a plateau; SD: one
% proton survives, a gap; DD: both dissociate, a gap in the middle; CD: both
% survive, a small central system), their cross sections and their shares in
% the 2000-event library of make_pileup_library.py (13.6 TeV).
% ---------------------------------------------------------------------------
\begin{tikzpicture}[x=1cm,y=1cm,
  every node/.style={font=\footnotesize},
  lab/.style={font=\scriptsize,align=center},
  head/.style={font=\small\bfseries},
  prot/.style={-{Stealth[length=2mm]},thick},
  prof/.style={thick}]
\foreach \dx/\name/\col/\code/\sig/\share/\np/\desc in {
   0/ND/blue!60!black/101/55.4 mb/71.2\%/230/{both protons break up; several parton pairs interact (MPI); a plateau of $\sim6$ charged per unit of $\eta$},
   5.2/SD/green!50!black/103+104/12.9 mb/15.9\%/51/{one proton survives intact, the other dissociates; a rapidity gap on the surviving side},
   10.4/DD/orange!80!black/105/8.8 mb/10.9\%/62/{both dissociate into low-mass systems; a gap in the middle, little in the tracker},
   15.6/CD/violet!80!black/106/1.3 mb/1.9\%/28/{both survive; a small system in the centre: a few particles}} {
 \begin{scope}[xshift=\dx cm]
  \node[head,anchor=west,text=\col] at (-0.2,3.3) {\name\quad\textcolor{black!60}{code \code}};
  % the eta axis
  \draw[-{Stealth}] (-0.2,0) -- (4.4,0) node[below left,lab] {$\eta$};
  \foreach \x/\t in {0.4/$-8$,1.3/$-4$,2.1/0,2.9/4,3.8/8} {\draw (\x,-0.05) -- (\x,0.05); \node[lab,anchor=north] at (\x,-0.08) {\t};}
  \draw[black!30] (1.6,0) -- (1.6,1.9); \draw[black!30] (2.6,0) -- (2.6,1.9);
  \node[lab,text=black!50] at (2.1,2.0) {tracker};
  \draw[black!15,dashed] (1.1,0) -- (1.1,1.9); \draw[black!15,dashed] (3.1,0) -- (3.1,1.9);
 \end{scope}
}
% the profiles, hand drawn: dN/deta per unit eta (scale 0.2 cm per particle)
\begin{scope}
 \draw[prof,blue!60!black] plot[smooth] coordinates {(0.3,0.1)(0.7,0.9)(1.1,1.3)(1.5,1.4)(2.1,1.45)(2.7,1.4)(3.1,1.3)(3.5,0.9)(3.9,0.1)};
 \node[lab,text=blue!60!black] at (2.1,1.7) {$\sim6$ charged per unit $\eta$};
\end{scope}
\begin{scope}[xshift=5.2cm]
 \draw[prof,green!50!black] plot[smooth] coordinates {(0.3,0.1)(0.7,0.8)(1.1,1.0)(1.4,0.9)(1.7,0.5)(1.9,0.15)(2.2,0.05)(3.9,0.05)};
 \draw[prot,green!50!black] (3.9,0.6) -- (4.3,0.6); \node[lab,text=green!50!black,anchor=south] at (4.0,0.65) {$p$};
 \node[lab,text=green!50!black] at (3.0,0.5) {gap};
\end{scope}
\begin{scope}[xshift=10.4cm]
 \draw[prof,orange!80!black] plot[smooth] coordinates {(0.3,0.1)(0.7,0.8)(1.0,0.9)(1.3,0.5)(1.6,0.1)(2.6,0.1)(2.9,0.5)(3.2,0.9)(3.5,0.8)(3.9,0.1)};
 \node[lab,text=orange!80!black] at (2.1,0.5) {gap};
\end{scope}
\begin{scope}[xshift=15.6cm]
 \draw[prof,violet!80!black] plot[smooth] coordinates {(0.3,0.03)(1.5,0.05)(1.9,0.4)(2.1,0.5)(2.3,0.4)(2.7,0.05)(3.9,0.03)};
 \draw[prot,violet!80!black] (3.9,0.6) -- (4.3,0.6); \draw[prot,violet!80!black] (0.3,0.6) -- (-0.1,0.6);
 \node[lab,text=violet!80!black,anchor=south] at (4.0,0.65) {$p$}; \node[lab,text=violet!80!black,anchor=south] at (0.2,0.65) {$p$};
\end{scope}
% the table below each
\foreach \dx/\col/\sig/\share/\np/\nc/\desc in {
   0/blue!60!black/55.4 mb/71.2\%/230/16/{both protons break up; several parton pairs interact (MPI); a plateau},
   5.2/green!50!black/12.9 mb/15.9\%/51/3/{one proton survives, the other dissociates; a gap on the surviving side},
   10.4/orange!80!black/8.8 mb/10.9\%/62/3/{both dissociate into low-mass systems; a gap in the middle, little in the tracker},
   15.6/violet!80!black/1.3 mb/1.9\%/28/2/{both survive; a small system in the centre, a few particles}} {
  \node[lab,anchor=north west,text width=4.5cm,align=left] at (\dx cm-0.2cm,-0.55) {\textcolor{\col}{$\sigma=\sig$, \share\ of the library}\\ \np\ particles on average, \nc\ charged in the tracker\\ \desc};
}
\node[lab,anchor=north west,text width=20cm,align=left] at (-0.2,-2.5) {\textbf{Minimum bias} is all of them in their natural proportions, $\sigma_{\rm inel}=78.4$ mb at 13.6 TeV (\PYTHIA, CP5: $\sigma_{\rm tot}=100.3$ and $\sigma_{\rm el}=21.9$ mb fixed by the tune, the same at 14 TeV; the diffractive pieces from its parametrisation, ND what remains). The library holds 2000 of them: 179 final particles per collision on average, a median of 116, from 3 to 829; the hardest particle $2.2$ GeV on average; $0.95$ GeV of $p_T$ per unit area in the tracker's reach.};
\end{tikzpicture}
}
\caption[One member of the crowd: the four kinds of inelastic collision]{One
member of the crowd. The four processes of \code{SoftQCD:inelastic} with
their profiles in rapidity, their cross sections at $13.6\TeV$, their
shares in the library of two thousand collisions and the particles each
brings on average.}
\label{fig:crowd_on_the_road:minbias_sketch}
\end{figure}

\Cref{tab:crowd_on_the_road:catalogue} is the crowd one member at a
time, the catalogue of the library as the companion's
\code{minBiasGenStats.py} prints it for the $250$ GEN events of the 2024
\CMSSW library: the first ten collisions of our in-time crossing, each
with its kind, its parton-parton interactions, its particles, what the
tracker and what nothing will see of them, its hardest particle, its
jets, its forward protons, and the vertex the mixing drew for it. The
first of the ten is a single-diffractive collision of four particles,
a surviving proton among them, that no central detector will notice; the
eighth is a non-diffractive one with $24$ interactions and $537$
particles, $93$ of them in the tracker's reach, more than the family's
own $93$. The summary below it, over the two thousand, and the means per
kind say the rest.

\begin{table}[htbp]\centering\small
\caption[The catalogue of the library: ten collisions, the summary and
the means per kind]{The catalogue of the library, written by
\file{minbias_catalogue.py} in the format of the companion's
\code{minBiasGenStats.py}. Top: the first ten in-time collisions of our
crossing, with their library event, kind, parton-parton interactions
(the hardest scattering plus the multiparton ones), stable particles,
charged ones, central ones ($|\eta|<2.5$ and $\pt>0.5\GeV$, all species),
the ones beyond $|\eta|=5.5$ where no subdetector is, the hardest
particle, the scalar sum of \pt, the \antikt $R=0.4$ jets above
$10\GeV$, the forward protons that keep most of the beam energy
($\xi<0.55$, the ones \CMSSW's \code{LHCTransport} carries to the PPS),
and the vertex the mixing drew. Middle: the same quantities over the two
thousand collisions of the library, with the vertices of the $960$
collisions of our event. Bottom: the means per kind, with the cross
section of each at $13.6\TeV$.}
\label{tab:crowd_on_the_road:catalogue}
\scriptsize\setlength{\tabcolsep}{3pt}
\resizebox{\textwidth}{!}{\minbiasCatalogueTen}\\[2mm]
\minbiasCatalogueSummary\\[2mm]
\resizebox{\textwidth}{!}{\minbiasCataloguePerProcess}
\end{table}

A member of the crowd, then, is soft and uneven. The library's
collisions have from $3$ to $829$ final particles, $179$ on average and a
median of $116$: a non-diffractive one brings $230$ particles and $16$
charged ones into the tracker's reach, $|\eta|<2.5$ and $\pt>0.5\GeV$, a
diffractive one $50$ to $60$ and $3$, a central-diffractive one $28$ and
$2$; the hardest particle of a collision is $2.2\GeV$ on average, and in
two thousand collisions it never exceeded $18$. Where they walk is the
plateau: $5.5$ charged particles per unit of pseudorapidity at $\eta=0$,
$13$ with the neutrals, falling off beyond the tracker and gone beyond
$|\eta|=8$, where the beam remnants carry the rest of the energy into the
beam pipe. What they carry is the number this chapter is about. The
transverse momentum of a collision per unit of $\eta$ and $\phi$,
\begin{equation}
  \rhoPU_1=\frac{1}{2\pi}\,\Big\langle\frac{\mathrm{d}\Sigma\pt}{\mathrm{d}\eta}\Big\rangle_{\text{one collision}}
  =0.95\GeV\ \text{per unit area at }|\eta|<2.5,
  \label{eq:crowd_on_the_road:rho1}
\end{equation}
averaged over the library without the neutrinos, is the slope of
everything that follows: sixty collisions put $57\GeV$ into every unit of
area, and a jet of $R=0.4$ covers $\pi R^2=0.503$ of them, $29\GeV$
under our $100\GeV$ jet before any subtraction.
\Cref{fig:crowd_on_the_road:theory_member} shows the three panels, and
the one surprise in them: the multiplicity distribution is so wide that
the median collision has two thirds of the mean, and in any crossing a
handful of the sixty, the non-diffractive ones with many multiparton
interactions, carry a quarter of the crowd. The \CMSSW library of the
companion, $250$ collisions through the full chain, has $194$ stable
particles per collision with a median of $132$, from $3$ to $723$, and
the same shape; ten of its events, single-diffractive ones that hand
\GEANT a surviving proton and nothing else, leave no hit in the central
detector at all.

\begin{figure}[htbp]\centering
\includegraphics[page=2,width=\textwidth]{PU/Standalone/theory_pileup.pdf}
\caption[One member of the crowd: the multiplicity, the plateau, the
transverse momentum per unit area]{Page~2 of \file{theory_pileup.pdf}.
(a) The final particles per collision in the library, by process type,
with the \CMSSW library's mean and our signal event for comparison. (b)
The particles per unit of $\eta$ per collision, all and charged: the
plateau, the tracker's reach and the beam remnants. (c) The transverse
momentum per unit area per collision, \cref{eq:crowd_on_the_road:rho1},
and the slope of \rhoPU it predicts against the one measured over four
hundred draws.}
\label{fig:crowd_on_the_road:theory_member}
\end{figure}

\subsection{Where and when: the beam spot and the 25 ns clock}
\label{sec:crowd_on_the_road:in_time_out_of_time_and_the_25_ns_bunch_}

The crowd walks in the same picture, but not from the same point. Two
bunches of protons pass through each other, and a collision can happen
anywhere in their overlap and at any moment of it; the region where the
collisions happen, the \emph{beam spot}, is what the vertex finder of
\cref{ch:first_checkpoint} will use to tell the crowd from the family, and
its time spread is what no detector of Run~3 can resolve.
\Cref{fig:crowd_on_the_road:region_sketch} draws it. Bunch~1 moves along
$+z$ and bunch~2 along $-z$; at the laboratory time $t$ their centres are
at $z=\pm ct$, and a proton at the longitudinal offset $s$ from its bunch
centre sits at $z=\pm ct+s$. For two Gaussian bunches of length
$\sigma_s$ meeting head-on the density of collisions is the product of
the two bunch densities at the same point and time,
\begin{equation}
  p(z,t)\propto
  \exp\!\Big[-\frac{(z-ct)^2+(z+ct)^2}{2\sigma_s^2}\Big]
  =\exp\!\Big[-\frac{z^2}{2\sigma_z^2}\Big]\exp\!\Big[-\frac{t^2}{2\sigma_t^2}\Big],
  \qquad \sigma_z=\frac{\sigma_s}{\sqrt2},\quad \sigma_t=\frac{\sigma_z}{c},
  \label{eq:crowd_on_the_road:region}
\end{equation}
after the transverse coordinates have been integrated out: a Gaussian
along the beam, a Gaussian in time of the same width in units of $z/c$,
and the two independent, because the cross term $z\,ct$ cancels between
the two bunches. Across the beam the collisions follow the beam size,
$\sigma_{x,y}(z)=\sqrt{\epsilon\beta^*/2}\,\sqrt{1+(z-z_0)^2/\beta^{*2}}$,
narrowest at the focus and growing on either side, the hourglass of a
beam focused to $\beta^*$. This is exactly the model of \CMSSW's
\code{BetafuncEvtVtxGenerator}, the \code{VtxSmeared} of every Run-3
GEN step: it draws $z$ and $t$ as independent Gaussians of the same
width, and $x$ and $y$ around the centre of the beam spot with the
hourglass width at that $z$. What it does not do is compute $\sigma_z$
from the bunches. The real bunches of Run~3 are longer than
$\sqrt2\times3.6\cm$ and cross at an angle, which shortens their overlap,
and the generator reads instead the \emph{measured} beam spot from the
database, the \code{SimBeamSpot} payload of the 2024 global tag, a fit to
the early-2024 data: centre $(0.03,-0.02,0.42)\cm$, $\sigma_z=3.6\cm$,
$\beta^*=30\cm$ and an emittance of $5.7\times10^{-8}\cm$, hence
$\sigma_t=120\ps$ and $\sigma_{xy}=9\um$ at the focus. The script draws
every vertex of the crowd, and the signal's, from that payload: the $960$
vertices of our event have an r.m.s.\ of $3.79\cm$ in $z$ and $119\ps$
in $t$, $9\um$ across, and the signal sits at $z=-0.17\cm$, $t=-121\ps$.
The $250$-event \CMSSW library of the companion, drawn by the real
module, has $3.56\cm$ and $122\ps$ with its vertices between $-9.5$ and
$+10.1\cm$, and the \CMSSW signal event sits at $z=3.11\cm$. The HL-LHC
of the open problems is another story: there the crossing angle would
cut the overlap to $2\cm$, crab cavities give most of it back, and
\CMSSW's \code{HLLHCEvtVtxGenerator} builds the region from the bunches,
$4.3\cm$ and $190\ps$; the script keeps that model as \code{--region
hllhc}.

\begin{figure}[htbp]\centering
\resizebox{\textwidth}{!}{% ---------------------------------------------------------------------------
% The beam spot of 2025: (a) the two bunches in the (x, z) plane and their
% overlap, with the hourglass of the beta function; (b) the numbers of the
% payload against the drawn vertices, the CMSSW library and the HL-LHC model;
% the clock of the sixteen crossings is pu_clock.tex.
% ---------------------------------------------------------------------------
\begin{tikzpicture}[x=1cm,y=1cm,
  every node/.style={font=\footnotesize},
  lab/.style={font=\scriptsize,align=center},
  head/.style={font=\small\bfseries}]
% ---------------- (a) the crossing plane --------------------------------------------------------------
\begin{scope}
\node[head,anchor=west] at (-0.3,3.4) {(a) the beam spot in the $(x,z)$ plane};
\draw[-{Stealth}] (-4.2,0) -- (4.4,0) node[below left,lab] {$z$};
\draw[-{Stealth}] (0,-1.6) -- (0,2.2) node[left,lab] {$x$ ($\times10^{3}$)};
% the hourglass envelope: sigma(z) = sigma0 sqrt(1 + z^2/beta*^2), beta* = 30 cm (scale 0.2 cm per cm)
\draw[black!35,dashed] plot[domain=-4:4,samples=60] (\x,{0.22*sqrt(1+(\x/1.2)*(\x/1.2))});
\draw[black!35,dashed] plot[domain=-4:4,samples=60] (\x,{-0.22*sqrt(1+(\x/1.2)*(\x/1.2))});
\node[lab,text=black!50,anchor=west] at (-4.1,1.7) {hourglass: $\sigma_x(z)=\sigma_{x0}\sqrt{1+(z-z_0)^2/\beta^{*2}}$, $\beta^*=30$ cm};
% bunch 1 (blue) moving +z, bunch 2 (orange) moving -z, a small crossing angle
\begin{scope}[rotate=5]
 \fill[blue!25,opacity=0.8] (-1.9,0) ellipse (1.8 and 0.22);
\end{scope}
\begin{scope}[rotate=-5]
 \fill[orange!35,opacity=0.8] (1.9,0) ellipse (1.8 and 0.22);
\end{scope}
\draw[blue!60!black,thick,-{Stealth[length=2mm]}] (-3.6,-0.75) -- (-2.6,-0.6) node[lab,below,pos=0.3] {bunch 1};
\draw[orange!80!black,thick,-{Stealth[length=2mm]}] (3.6,-0.75) -- (2.6,-0.6) node[lab,below,pos=0.3] {bunch 2};
\fill[green!50!black,opacity=0.55] (0.1,0) ellipse (0.72 and 0.16);
\node[lab,text=green!50!black,anchor=north] at (0.1,-0.25) {the overlap: $\sigma_z=3.6$ cm, $\sigma_t=120$ ps, $\sigma_{xy}=9\,\mu$m};
\node[lab,anchor=north west,text width=8.8cm,align=left] at (-4.2,-1.85) {Two Gaussian bunches meeting head-on collide over $\sigma_z=\sigma_s/\sqrt2$, during $\sigma_t=\sigma_z/c$, with $z$ and $t$ independent, \cref{eq:crowd_on_the_road:region}. The crossing angle of Run 3 shortens the overlap; \CMSSW does not model it but reads the measured region from the database: \code{BetafuncEvtVtxGenerator} draws $z$, $t$, $x$ and $y$ from the 2025 payload, the signal and every library event alike.};
\end{scope}
% ---------------- (b) the numbers --------------------------------------------------------------------
\begin{scope}[xshift=12.4cm]
\node[head,anchor=west] at (-3.2,3.4) {(b) how long, how wide: the payload and three checks};
\foreach \y/\t/\v/\w in {2.6/{the 2025 payload}/3.60/2.88, 1.9/{our 960 vertices}/3.56/2.85, 1.2/{\CMSSW library (1000)}/3.57/2.86, 0.5/{the HL-LHC model}/4.27/3.42} {
  \draw[fill=green!30] (0,\y-0.18) rectangle (\w,\y+0.18);
  \node[lab,anchor=west] at (\w+0.1,\y) {\v\ cm};
  \node[lab,anchor=east] at (-0.1,\y) {\t};
}
\draw[-{Stealth}] (0,0.1) -- (4.3,0.1) node[below left,lab] {$\sigma_z$ [cm], 0.8 cm per cm};
\node[lab,anchor=north west,text width=7.6cm,align=left] at (-3.2,-0.5) {In time the same four read $120$, $124$, $122$ and $193$ ps. The payload of the 2025 global tag is a fit to early-2024 data: centre $(0.03,-0.02,0.42)$ cm, $\sigma_z=3.6$ cm, $\beta^*=30$ cm, emittance $5.7\times10^{-8}$ cm, hence $\sigma_{xy}=\sqrt{\epsilon\beta^*/2}=9\,\mu$m at the focus. The HL-LHC region is longer because the crab cavities give back what its larger crossing angle takes.};
\end{scope}
\end{tikzpicture}
}
\caption[The beam spot of 2024: the two bunches and the numbers of the
payload]{The beam spot. (a) The $(x,z)$ plane: the two bunches, their
overlap and the hourglass of the $\beta$ function. (b) The length of the
beam spot: the 2024 payload against the vertices of our event, the
\CMSSW library and the HL-LHC model.}
\label{fig:crowd_on_the_road:region_sketch}
\end{figure}

The clock is the other half, and \cref{fig:crowd_on_the_road:clock}
draws it with the numbers of our event. The crossings before and after
ours, $25\ns$ apart, collide on the same beam spot, and their collisions
are the same crowd shifted in time: the \Mix draws each of the sixteen
crossings of the production window on its own, $60$, $56$, $68$, $61$,
$66$, $56$, $55$, $50$, $53$, $66$, $61$ and $60$ collisions in the
twelve early ones, $63$ in ours, $63$, $59$ and $63$ in the three late
ones, $960$ collisions piled on one signal event, of which $712$ arrive
early, up to $300\ns$ before the signal, and $185$ late, up to $75\ns$
after it. \CMSSW's \code{Adjuster} stamps every \code{SimTrack},
\code{SimVertex} and hit of collision $k$ of crossing \bx with
\code{EncodedEventId(bx, k)} and adds $25\ns\times\bx$ to its time, so
that in the simulation, and only there, the question ``which crowd did
this come from'' has an answer: \code{bunchCrossing() = 0} with
\code{event() = 0} is the signal, \code{bunchCrossing() = 0} with
\code{event()} of one or more is an in-time collision, and anything with
\code{bunchCrossing()} different from zero is out of time; the
\code{TrackingParticle}s and the \code{DigiSimLink}s carry the id
through the digitisation, which is how the companion's pages colour
every hit by its crowd and how a reconstructed track can be called
signal, in-time or out-of-time pileup in a validation. The record of
this stop keeps the same two numbers, \bx and $k$, in its vertex listing
for the $897$ out-of-time collisions and for the $63$ in-time ones. In
space nothing distinguishes an early collision from an in-time one, which
is why the tracker, which does not keep time to a nanosecond, cannot; the
calorimeter pulses of \cref{ch:first_checkpoint}, which last several
crossings, will integrate them, and the price of that is the out-of-time
pileup of the chapters ahead. At the generator level nothing integrates
over $25\ns$, so the particles of the other crossings are not clustered
here; they are in the record for \cref{ch:crossing_the_border}. What the
in-time crowd does to
the vertex finder, on the other hand, can be counted now. A collision the
tracker can see needs charged particles above its threshold pointing back
to a common point, and with the proxy of at least four charged particles
of $\pt>0.5\GeV$ within $|\eta|<2.5$, $43$ of the $63$ collisions of our
crossing are visible, $62\%$ over the ensemble, against the $0.7\mup$ that
CMS reconstructs with its real threshold; the rest are diffractive or
forward and leave the tracker empty. Their vertices are a Gaussian of
$3.6\cm$ with sixty entries, so the typical spacing is a millimetre: the
nearest crowd vertex to ours is $0.06\,\mathrm{mm}$ away in $z$ in our
crossing and $0.98$ on the median of the ensemble, in one draw of two one
of them is within a millimetre, and the vertex finder, which resolves
vertices a few hundred microns apart, will merge the two in a fair
fraction of the events; the \CMSSW event of the companion has three
pileup vertices within a millimetre of the signal's. In time the nearest
is five picoseconds away and nine of them are within $30\ps$ on average,
which is the number a timing layer would sort by; Run~3 has none, and in
this book the time of a collision is a label the record carries and no
detector reads. \Cref{fig:crowd_on_the_road:theory_region} draws the
density, the length of the beam spot three ways and the clock.

\begin{figure}[htbp]\centering
\resizebox{\textwidth}{!}{% ---------------------------------------------------------------------------
% The clock: the sixteen crossings of the Run-3 production window, 25 ns
% apart, with the number of collisions the MixingModule drew for each in our
% event (the -v log, seed 1), the in-time crossing with the signal, and the
% EncodedEventId(bx, k) that lets the simulation tell them apart.
% ---------------------------------------------------------------------------
\begin{tikzpicture}[x=1cm,y=1cm,
  every node/.style={font=\footnotesize},
  lab/.style={font=\scriptsize,align=center},
  head/.style={font=\small\bfseries}]
\node[head,anchor=west] at (-0.6,3.6) {the window $bx=-12\ldots+3$: sixteen crossings, sixteen Poisson draws, 960 collisions on one signal event};
\draw[-{Stealth}] (-0.6,0) -- (17.2,0) node[below left,lab] {$t$ [ns]};
% bars: height 0.04 cm per collision (60 -> 2.4 cm)
\foreach \bx/\n/\col in {-12/60/yellow!60!brown,-11/56/yellow!60!brown,-10/68/yellow!60!brown,-9/61/yellow!60!brown,-8/66/yellow!60!brown,-7/56/yellow!60!brown,-6/55/yellow!60!brown,-5/50/yellow!60!brown,-4/53/yellow!60!brown,-3/66/yellow!60!brown,-2/61/yellow!60!brown,-1/60/yellow!60!brown,0/63/cyan!60!blue,1/63/red!50!orange,2/59/red!50!orange,3/63/red!50!orange} {
  \pgfmathsetmacro{\x}{12.6+1.0*\bx}
  \pgfmathsetmacro{\h}{0.04*\n}
  \fill[\col,opacity=0.75] (\x-0.32,0) rectangle (\x+0.32,\h);
  \node[lab,anchor=south] at (\x,\h+0.02) {\n};
  \node[lab,anchor=north] at (\x,-0.06) {$\bx$};
}
\foreach \bx in {-12,-8,-4,0,3} {\pgfmathsetmacro{\x}{12.6+1.0*\bx} \node[lab,anchor=north,text=black!60] at (\x,-0.42) {$25\times(\bx)$};}
\draw[black!40,dashed] (-0.6,2.4) -- (16.4,2.4) node[lab,anchor=west,text=black!60] at (16.45,2.4) {$\mu=60$};
\node[star,star points=5,fill=red!80!black,inner sep=0pt,minimum size=9pt] at (12.6,1.2) {};
\node[lab,text=red!80!black,anchor=west] at (12.85,1.25) {ours};
\node[lab,anchor=north west,text width=5.2cm,align=left,text=yellow!50!brown] at (-0.6,-0.95) {\textbf{early}, $bx<0$: 12 crossings, 712 collisions; their hits arrive $25\,{\rm ns}\times|bx|$ before the signal, up to 300 ns, and the slow detectors (calorimeter pulses, strip readout) remember them};
\node[lab,anchor=north west,text width=4.6cm,align=left,text=cyan!50!blue] at (5.4,-0.95) {\textbf{in time}, $bx=0$: 63 collisions in the same crossing as the signal, 9033 particles in the detector's reach; the only ones a generator-level jet can see};
\node[lab,anchor=north west,text width=4.6cm,align=left,text=red!50!orange] at (10.6,-0.95) {\textbf{late}, $bx>0$: 3 crossings, 185 collisions, up to 75 ns after the signal: the tails of the latest pulses};
\node[lab,anchor=north west,text width=17.6cm,align=left] at (-0.6,-2.75) {In space the sixteen crowds sit on the same beam spot; only the clock tells them apart, and in the simulation the clock is written down: the \code{Adjuster} of the \code{MixingModule} stamps every \code{SimTrack}, \code{SimVertex} and hit of collision $k$ of crossing $bx$ with \code{EncodedEventId(bx, k)} and adds $25\,{\rm ns}\times bx$ to its time. \code{bunchCrossing() = 0} and \code{event() = 0} is the signal; \code{bunchCrossing() = 0} and \code{event() $\ge$ 1} an in-time pileup collision; \code{bunchCrossing() $\neq$ 0} an out-of-time one. The record of this stop keeps the same two numbers in its vertex listing ($bx$ and $k$), and the reconstruction of \cref{ch:first_checkpoint} will never see them: the truth of which hit came from which crowd exists in the simulation and nowhere else.};
\end{tikzpicture}
}
\caption[The clock: the sixteen crossings of the production window and
the collisions of each]{The clock. The sixteen crossings of the Run-3
production window, $25\ns$ apart, with the number of collisions the
\Mix drew for each in our event: $712$ early, $63$ in time with the
signal, $185$ late, $960$ in all; and the \code{EncodedEventId} by which
the simulation, and nothing after it, knows which crowd a hit came
from.}
\label{fig:crowd_on_the_road:clock}
\end{figure}

\begin{figure}[htbp]\centering
\includegraphics[page=3,width=\textwidth]{PU/Standalone/theory_pileup.pdf}
\caption[The beam spot: the density in $(z,t)$, the $z$ of the vertices,
the sixteen crossings]{Page~3 of \file{theory_pileup.pdf}. (a) The
collision density of \cref{eq:crowd_on_the_road:region} in $(z,t)$ with
the $63$ in-time collisions of our crossing and the signal. (b) The $z$
of the $960$ vertices of the event against the Gaussian of the payload
and the $250$ vertices of the \CMSSW library. (c) The sixteen crossings
on the clock, $25\ns$ apart, early in dark yellow, in time in blue and
late in salmon, the colours of the \CMSSW pages of the companion.}
\label{fig:crowd_on_the_road:theory_region}
\end{figure}

\paragraph{How \CMSSW does it: the library, the mixing and the premixing.}
\label{sec:crowd_on_the_road:how_cms_mixes_pileup_into_simulation_the}
What the script does with four-vectors, \CMSSW does with the hits of
\GEANT, and the strategy is the same. \GEANT is the expensive step,
$2.5$ seconds per minimum-bias collision, and the hits it leaves add: a
sensor crossed by particles of two collisions holds two hits, a
calorimeter cell the sum of two energies. So the minimum-bias collisions
are generated and simulated once, with the same fragment, the same vertex
smearing, the same geometry and the same \GEANT as the signal, but with
their own random seeds, and stored in a GEN-SIM library; the \Mix of the
DIGI step draws the number of collisions per crossing, reads that many
library events, stamps and shifts them, and hands the signal and the
pileup hits together to the digitisers, which add them channel by
channel and apply the pulse shapes, the time windows and the thresholds
in which the out-of-time pileup does its damage. The truth survives as
the \code{PileupSummaryInfo}, the \code{TrackingParticle}s with their
event ids, and the \code{DigiSimLink}s that say which collision
contributed what to each channel. The companion's exercise does exactly
this for one signal event, $\mup=60$ in three crossings, $66$, $55$ and
$77$ collisions drawn in sequence from a library of $250$: the $198$
collisions bring $31\,597$ primaries against the $571$ of the signal,
$444\,535$ tracker hits against $9776$, and own $104\,277$ of the
$106\,581$ hit cells of the ECAL. In the production window of sixteen
crossings this classical mixing would read a thousand library events, a
quarter of a gigabyte, per signal event, and the production therefore
\emph{premixes}~\cite{Hildreth:2017caz}: the minimum-bias events are
mixed among themselves at the flat profile and digitised once into a
library of premixed crossings, and a single premixed event is added, in
ADC counts, to every digitised signal event. The physics is the same,
the bookkeeping is a thousand times lighter, and the price is that the
pileup hits are gone from a premixed sample, so that pictures like those
of the companion, which colour every hit by the crossing it came from,
can only be drawn with the classical mixing; the script reproduces the
conditions of the exercise at the generator level with \code{--window 1
--sequential}.

\subsection{\rhoPU, the median energy density, and the jet area $A$}
\label{sec:crowd_on_the_road:rho_the_median_energy_density_and_the_je}

The crowd is uniform on average, and that is the handle. If every unit of
area of the $(\eta,\phi)$ plane receives \rhoPU GeV of transverse momentum
from the crowd, then a jet of area $A$ receives $\rhoPU A$, and
subtracting it is a matter of measuring \rhoPU in the parts of the plane
where nothing else is and $A$ for the jet. Both ideas are Cacciari and
Salam's~\cite{Cacciari:2007fd}, and \cref{fig:crowd_on_the_road:rho_sketch}
draws them. The density is estimated as a median: cut the plane into
patches, measure the transverse momentum per unit area in each, and take
the middle value,
\begin{equation}
  \rhoPU=\operatorname*{median}_{\text{patches }i}\Big[\frac{\pt^{(i)}}{A^{(i)}}\Big],
  \label{eq:crowd_on_the_road:rho}
\end{equation}
because the median does not care about the tail. The patches that hold a
jet are the top few percent and would pull a mean up, $74$ against $68$
in our crossing, and the median ignores them as long as fewer than half
the patches are jets, which is true of any event. Two choices of patch are
in use. CMS's \code{fixedGridRhoFastjetAll}, the \rhoPU of the jet energy
corrections of \cref{ch:paying_for_uninvited}, is \FASTJET's
\code{GridMedianBackgroundEstimator} on the particles of the whole
detector: a fixed grid of $18$ by $11$ cells of $0.56\times0.57$ over
$|y|<5$, in rapidity and with $\phi$ counted from zero, as \FASTJET draws
it; the original of \cite{Cacciari:2007fd} uses the $k_T$ jets of
$R=0.4$ as patches, each with its own area, within $|y|<4.4$, since the
$k_T$ algorithm of \cref{ch:the_inn} tiles the plane with jets of
comparable area rather than finding hard cores. The script does both. The
grid gives $\rhoPU=68.28\GeV$ per unit area in our crossing, the cells
ranging from $13$ to $359$; the $k_T$ jets, $175$ of them within the
range, give $69.87$; the signal alone, the underlying event of
\cref{ch:uninvited_cousins} and nothing in most cells, gives $3.42$. Per
collision that is $1.08$ and $1.11\GeV$ per unit area in this crossing,
a crowded one, and over four hundred draws \rhoPU is $59\pm12$ with a
slope of $0.94\GeV$ per collision and an offset of $3$, the family's own
underlying event, which the median counts as crowd: the slope is the
$\rhoPU_1$ of \cref{eq:crowd_on_the_road:rho1} within the difference
between a median and a mean, the $8$ collisions of Poisson spread
account for $7\GeV$ of the $12$, and the rest is the crowd's own
unevenness, the same number of collisions bringing more or fewer
particles.

\begin{figure}[htbp]\centering
\resizebox{\textwidth}{!}{% ---------------------------------------------------------------------------
% rho and the area: (a) the (eta, phi) plane tiled in the 18 x 11 cells of
% fixedGridRhoFastjetAll, the pT per unit area of each cell as a grey level,
% the jets dark; (b) the cells sorted: the median is rho, the mean is not;
% (c) the active area of a jet: the ghosts it swallowed times 0.01.
% ---------------------------------------------------------------------------
\begin{tikzpicture}[x=1cm,y=1cm,
  every node/.style={font=\footnotesize},
  lab/.style={font=\scriptsize,align=center},
  head/.style={font=\small\bfseries},
  ghost/.style={circle,fill=black!35,inner sep=0pt,minimum size=1.8pt},
  fam/.style={circle,fill=red!80!black,inner sep=0pt,minimum size=3.6pt},
  crowd/.style={circle,fill=olive!60!brown,inner sep=0pt,minimum size=2.4pt}]
% ---------------- (a) the grid ---------------------------------------------------------------------
\begin{scope}
\node[head,anchor=west] at (-0.3,3.9) {(a) the grid: 18 $\times$ 11 cells of $0.56\times0.57$ in $|y|<5$};
\pgfmathsetseed{5}
\foreach \i in {0,...,17} {
  \foreach \j in {0,...,10} {
    \pgfmathsetmacro{\v}{min(100, 35 + 25*rnd + 15*rnd*rnd)}
    \pgfmathsetmacro{\g}{100-\v*0.6}
    \fill[black!\g] (\i*0.4,\j*0.3) rectangle (\i*0.4+0.4,\j*0.3+0.3);
  }
}
\fill[black!85] (3.2,0.6) rectangle (3.6,0.9); \fill[black!75] (2.8,0.6) rectangle (3.2,0.9);
\fill[black!85] (2.0,2.4) rectangle (2.4,2.7); \fill[black!70] (4.8,1.5) rectangle (5.2,1.8);
\draw[black!40,very thin] (0,0) grid[step=0.4] (7.2,3.3);
\draw[red!80!black,thick] (3.3,0.75) circle (0.29);
\node[lab,text=red!80!black,anchor=west] at (3.65,0.75) {our jet};
\draw[-{Stealth}] (0,-0.25) -- (7.2,-0.25) node[below left,lab] {$y$: $-5$ to $5$};
\draw[-{Stealth}] (-0.25,0) -- (-0.25,3.3) node[above,lab] {$\phi$};
\node[lab,anchor=north west,text width=7.4cm,align=left] at (-0.3,-0.65) {Every particle's $p_T$ goes to its cell; the cell's value is $\Sigma p_T/(\Delta y\,\Delta\phi)$. At $\mu=60$ a typical cell holds $\sim22$ GeV in $0.32$ units of area, 68 GeV per unit area; the cells with a jet hold up to 360.};
\end{scope}
% ---------------- (b) the median --------------------------------------------------------------------
\begin{scope}[xshift=8.6cm]
\node[head,anchor=west] at (-0.4,3.9) {(b) the 198 cells sorted: the median is $\rho$};
\draw[-{Stealth}] (0,0) -- (6.0,0) node[below left,lab] {cell, sorted by $p_T$/area};
\draw[-{Stealth}] (0,0) -- (0,3.4) node[left,lab] {GeV per unit area};
% a sorted distribution: flat-ish middle, a tail of jets at the right
\draw[olive!60!brown,very thick] plot[domain=0:5.7,samples=80] (\x,{0.25+1.2*\x/5.7+0.25*(\x/5.7)^3+2.2*exp(-(5.7-\x)*(5.7-\x)/0.25)});
\draw[red!80!black,thick,dashed] (2.85,0) -- (2.85,2.6);
\node[lab,text=red!80!black,anchor=south] at (2.85,2.65) {median: $\rho=68$};
\draw[black!60,dashed] (0,1.25) -- (5.7,1.25);
\node[lab,text=black!60,anchor=west] at (0.1,1.45) {mean, pulled up by the tail};
\node[lab,text=olive!60!brown,anchor=east] at (5.6,3.2) {the jets};
\node[lab,anchor=north west,text width=6.4cm,align=left] at (-0.4,-0.65) {$\rho=\mathrm{median}_{\rm cells}\,[\,p_T/A\,]$. The median is what the crowd puts in a cell with nothing else in it; the jets are the top few percent and do not move it. Cacciari and Salam's original uses $k_T$ jets of $R=0.4$ with their areas in place of the cells: $\rho_{\rm jets}$ a few percent above the grid ($69.9$ against $68.3$ here).};
\end{scope}
% ---------------- (c) the area ------------------------------------------------------------------------
\begin{scope}[xshift=16.2cm]
\node[head,anchor=west] at (-0.4,3.9) {(c) the active area: ghosts of $0.01$ each};
\pgfmathsetseed{3}
\foreach \i in {0,...,15} {\foreach \j in {0,...,13} {
  \pgfmathsetmacro{\xx}{0.1+0.33*\i+0.2*(rnd-0.5)}
  \pgfmathsetmacro{\yy}{0.1+0.235*\j+0.15*(rnd-0.5)}
  \node[ghost] at (\xx,\yy) {};
}}
\draw[red!80!black,thick] (2.7,1.65) circle (1.15);
\pgfmathsetseed{11}
\foreach \i in {1,...,22} {\pgfmathsetmacro{\r}{1.15*0.5*rnd*rnd} \pgfmathsetmacro{\a}{360*rnd} \node[fam] at ({2.7+\r*cos(\a)},{1.65+\r*sin(\a)}) {};}
\pgfmathsetseed{17}
\foreach \i in {1,...,40} {\pgfmathsetmacro{\xx}{0.2+5.0*rnd} \pgfmathsetmacro{\yy}{0.15+3.0*rnd} \node[crowd] at (\xx,\yy) {};}
\node[lab,anchor=north west,text width=6.0cm,align=left] at (-0.4,-0.1) {A grid of ghosts, $p_T=10^{-18}$, one per $0.01$ of the plane, is clustered with the particles. They move nothing (anti-$k_T$ never sees their $p_T$) but they are caught like any soft particle: $A=0.01\times n_{\rm ghosts}$ in the jet. For a lone hard jet $A\simeq\pi R^2=0.503$; a neighbour steals part of it. Then $p_T^{\rm sub}=p_T-\rho\,A$: the \code{L1FastJet} correction of \cref{ch:paying_for_uninvited} in its generator-level form.};
\end{scope}
\end{tikzpicture}
}
\caption[\rhoPU and the area: the grid, the median, the ghosts]{\rhoPU
and the area. (a) The $(y,\phi)$ plane tiled in the cells of
\code{fixedGridRhoFastjetAll}, the transverse momentum per unit area of
each cell as a grey level, the jets dark. (b) The cells sorted: the
median is \rhoPU, the mean is pulled up by the tail. (c) The active area
of a jet: the ghosts it swallowed, $0.01$ of the plane each, and the
subtraction $\pt-\rhoPU A$.}
\label{fig:crowd_on_the_road:rho_sketch}
\end{figure}

The area is measured with the trick of \cref{ch:family_photograph}. A
grid of ghosts, one per $0.01$ of the plane within $|y|<5$, $6300$ of
them scattered within their cells with a transverse momentum of
$10^{-18}\GeV$, is clustered together with the particles; \antikt never
sees their momentum, since the distance to a soft particle is set by the
hard one, but it catches them like any soft particle within $R$ of the
core, and the \emph{active area} of a jet is the number of ghosts it
holds times the area of one,
\begin{equation}
  A_j=0.01\times n_{\mathrm{ghosts}}(j),\qquad
  \pt^{\mathrm{sub}}=\pt-\rhoPU\,A_j,
  \label{eq:crowd_on_the_road:area}
\end{equation}
\FASTJET's \code{active\_area\_explicit\_ghosts}~\cite{Cacciari:2008gn} with
CMS's \code{GhostArea = 0.01}. A lone hard jet has $A\simeq\pi R^2=0.503$,
and our jet, whose neighbour at $\dR=0.56$ no longer competes for the
soft particles between them once the crowd fills the gap, has $0.529$;
the leading twelve jets of the crossing range from $0.49$ to $0.56$. The
subtraction is then one line, \cref{eq:crowd_on_the_road:area}, the
\code{L1FastJet} correction of \cref{ch:paying_for_uninvited} in its
generator-level form, and the script can check it against a truth the
detector never has, because every constituent of a mixed jet carries the
collision it came from: our jet reads $132.75\GeV$ with the crowd,
$\rhoPU A=68.28\times0.529=36.09$, $\pt^{\mathrm{sub}}=96.66$, and the
family's particles inside it sum to $100.07$; the crowd had put $33.61$
in, from $81$ particles of $34$ different collisions, and the subtraction
took $36.09$ out, a residual of $-3.41\GeV$. That is one draw, and the
ensemble says how typical: over four hundred draws the crowd puts
$27.8\pm8.4\GeV$ into the circle of our jet, $\rhoPU\pi R^2$ estimates
$29.7\pm5.8$, and the residual is $-1.9\pm6.7\GeV$. Two lessons are in
those three numbers. The subtraction removes the mean, and a little
more: thirty GeV of strangers become minus two, because the median
measured over the whole plane includes the family's own underlying event,
$3\GeV$ per unit area, which is signal in the truth and crowd to the
estimator; the offset corrections of \cref{ch:paying_for_uninvited}
measure and remove exactly this kind of bias. And it cannot remove the
fluctuation: the crowd inside a circle of area $0.5$ varies by $\pm8\GeV$
from crossing to crossing, the estimate $\rhoPU A$ follows only the part
of that which is common to the whole plane, and $\pm7\GeV$ remain, on a
$100\GeV$ jet whose own photograph was stable to half a GeV. This is the
first term of the jet resolution of \cref{ch:wobbly_ruler} in Run~3, and
the reason the chapters between here and there exist: a subtraction that
knows which particles came from which vertex, \cref{ch:the_police}, can
do what a density cannot. \Cref{fig:crowd_on_the_road:theory_rho} shows
the cells with their median, \rhoPU against the number of collisions, and
the truth against the estimate over the ensemble.

\begin{figure}[htbp]\centering
\includegraphics[page=4,width=\textwidth]{PU/Standalone/theory_pileup.pdf}
\caption[\rhoPU and the area: the cells and their median, \rhoPU against
the crowd, the truth against $\rhoPU A$]{Page~4 of \file{theory_pileup.pdf}.
(a) The $198$ grid cells of our crossing, their median \rhoPU and their
mean. (b) \rhoPU against the number of in-time collisions over four
hundred draws, with the slope per collision and our crossing. (c) The
pileup inside the circle of our jet against the estimate $\rhoPU\pi R^2$,
draw by draw, with the residual in the inset.}
\label{fig:crowd_on_the_road:theory_rho}
\end{figure}

\paragraph{A word on the clustering.}
Sixteen thousand pseudojets, particles and ghosts, are ten times what
\cref{ch:family_photograph} clustered, and the plain search of its
\code{Clustering} class, which revisits every pseudojet whose nearest
neighbour has just merged, slows down badly under \antikt, where the
nearest neighbour of every ghost is a hard particle, so that each merge
of a hard particle touches hundreds of ghosts. The script adds the one
idea of \FASTJET that makes the problem affordable, the tiled search: a
pair farther apart than $R$ is never merged, because $d_{ij}$ is then
larger than the beam distance of the softer of the two, which becomes a
jet first; so the nearest neighbour of a pseudojet lies within $R$ of it,
the plane is tiled in cells of side $R$, and both the search and the
update look only at the nine tiles around a pseudojet. The result is the
same, step by step; the $15\,984$ steps of our crossing take two minutes
instead of an hour, and the $33\,000$ of an HL-LHC crossing a quarter of
an hour instead of a day.

\subsubsection{The recipe, step by step}
\label{sec:crowd_on_the_road:recipe}
\Cref{fig:crowd_on_the_road:flow} is the script as a flow chart, with the
numbers of our event beside each step: the scenario and \mup, the sixteen
Poisson draws and the library, the beam spot and the vertices, the
stacking of the in-time crowd, the tiled \antikt with ghosts, the two
medians, the subtraction against the truth, the visible vertices and the
profile, and the record.

\begin{figure}[htbp]\centering
\resizebox{\textwidth}{!}{% ---------------------------------------------------------------------------
% The stage as a flow chart: the steps of standalone_pileup.py on the left,
% the numbers of our event (the -v log, seed 1, mu = 60, bx = -12 ... +3) on the right. Names
% of the classes and functions in grey.
% ---------------------------------------------------------------------------
\begin{tikzpicture}[x=1cm,y=1cm,
  every node/.style={font=\footnotesize},
  box/.style={draw,rounded corners=3pt,thick,fill=gray!8,align=left,inner sep=5pt,minimum height=10mm,text width=5.6cm},
  stepbox/.style={box,fill=orange!12},
  numbox/.style={draw,rounded corners=3pt,thin,fill=blue!6,align=left,inner sep=4pt,text width=8.6cm,font=\scriptsize},
  stepno/.style={draw,circle,inner sep=1pt,font=\scriptsize\bfseries,fill=orange!45,minimum size=5mm},
  flow/.style={-{Stealth[length=2.5mm,width=2mm]},thick},
  fn/.style={font=\scriptsize\ttfamily,text=black!65},
  head/.style={font=\small\bfseries}]
\node[head,anchor=west] at (-0.6,1.4) {the recipe};
\node[head,anchor=west] at (6.7,1.4) {our event (seed 1, $\mu=60$, $bx=-12\ldots+3$)};
\node[stepbox] (S1) at (2.5,0) {\textbf{1\; the scenario}\\ $\mu=L\sigma_{\rm inel}/(n_b f_{\rm rev})$ from the luminosity; the window of bunch crossings, 25 ns apart};
\node[stepbox,below=7mm of S1] (S2) {\textbf{2\; the draw}\\ for every crossing $n\sim{\rm Poisson}(\mu)$; the library read in order from a random position};
\node[stepbox,below=7mm of S2] (S3) {\textbf{3\; the beam spot}\\ $(x,y,z,t)$ of every collision, the signal included, from the 2025 payload: Gaussian in $z$ and $t$, the hourglass across; $+25\,{\rm ns}\times bx$};
\node[stepbox,below=7mm of S3] (S4) {\textbf{4\; the stacking}\\ one final state: the signal and the in-time collisions, $|\eta|<5$, no neutrinos; the out-of-time ones recorded, not clustered};
\node[stepbox,below=7mm of S4] (S5) {\textbf{5\; the clustering}\\ anti-$k_T$, $R=0.4$, with a grid of ghosts of area 0.01 (tiled search); every constituent tagged by its collision};
\node[stepbox,below=7mm of S5] (S6) {\textbf{6\; $\rho$ and $A$}\\ the grid median of the cells ($0.56\times0.57$, $|y|<5$) and the $k_T$-jet median ($|y|<4.4$); $A=0.01\,n_{\rm ghosts}$};
\node[stepbox,below=7mm of S6] (S7) {\textbf{7\; the subtraction}\\ $p_T^{\rm sub}=p_T-\rho A$ for every jet, against the truth: the signal $p_T$ inside it};
\node[stepbox,below=7mm of S7] (S8) {\textbf{8\; the vertices, the profile}\\ the collisions the tracker can see; the flat profile of the production, the profile of a year and the weight $P_{\rm target}/P_{\rm sample}$};
\node[stepbox,below=7mm of S8] (S9) {\textbf{9\; write}\\ the record of chapter 7, then the vertex, particle, jet and constituent tables of this stop};
\foreach \a/\b in {S1/S2,S2/S3,S3/S4,S4/S5,S5/S6,S6/S7,S7/S8,S8/S9} {\draw[flow] (\a) -- (\b);}
\foreach \n/\s in {1/S1,2/S2,3/S3,4/S4,5/S5,6/S6,7/S7,8/S8,9/S9} {\node[stepno] at ($(\s.west)+(-0.35,0)$) {\n};}
\node[fn,anchor=south east] at (S1.north east) {mu\_of, SCENARIOS};
\node[fn,anchor=south east] at (S2.north east) {Stage.draw\_crowd, Library};
\node[fn,anchor=south east] at (S3.north east) {BeamSpot.draw (LuminousRegion: HL-LHC)};
\node[fn,anchor=south east] at (S4.north east) {Stage.stack};
\node[fn,anchor=south east] at (S5.north east) {TiledClustering, MixedJets, ghost\_grid};
\node[fn,anchor=south east] at (S6.north east) {grid\_rho, jet\_median\_rho};
\node[fn,anchor=south east] at (S7.north east) {Stage.subtract};
\node[fn,anchor=south east] at (S8.north east) {visible\_vertices, fill\_profile, reweight};
\node[fn,anchor=south east] at (S9.north east) {pileup\_tables, read\_pileup\_tables};
\node[numbox,anchor=west] (N1) at (6.7,0) {$2.23\times10^{34}$ cm$^{-2}$s$^{-1}$, 78.4 mb, 2340 bunches, 11\,245 Hz: $\mu=66.5$; mixed at 60, window $bx=-12\ldots+3$, 960 collisions expected (the production: $\mu_{\rm true}$ flat in 55--75, $\sim1040$)};
\node[numbox,anchor=west] (N2) at (6.7,0 |- S2) {sixteen draws from 50 to 68, \textbf{63} in time, 960 in all; library events 430 to 1389, each once (ND 698, SD 158, DD 86, CD 18)};
\node[numbox,anchor=west] (N3) at (6.7,0 |- S3) {payload: $\sigma_z=3.6$ cm, $\sigma_t=120$ ps; drawn: r.m.s.\ 3.56 cm, 124 ps, $x$ and $y$ 9 $\mu$m; the signal at $z=+0.16$ cm, $t=+94$ ps; nearest neighbour 0.01 mm and 0.3 ps away};
\node[numbox,anchor=west] (N4) at (6.7,0 |- S4) {signal: 651 of 744 particles; in time: 63 collisions, 13\,456 particles, 10\,045 within $|\eta|<5$; 995 charged in the tracker against 93 of the signal};
\node[numbox,anchor=west] (N5) at (6.7,0 |- S5) {10\,696 particles + 6300 ghosts, 16\,996 steps (3 min); 175 jets above 3 GeV, 82 above 30; our jet, mixed jet [1]: 115 particles, 31 of the family and 84 strangers from 33 collisions, $A=0.529$};
\node[numbox,anchor=west] (N6) at (6.7,0 |- S6) {grid: $\rho=75.5$ (cells from 30 to 361, mean 83); $k_T$ jets: 78.5; per collision 1.20 and 1.25 GeV per unit area in this crossing, 0.97 over the ensemble; the signal alone: 3.4, the crowd alone: 69.9};
\node[numbox,anchor=west] (N7) at (6.7,0 |- S7) {our jet $133.53-75.47\times0.529=93.64$ against $100.07$ of signal inside: residual $-6.43$ GeV; the crowd had put $34.57$ in; 76 pileup jets above 30 GeV before, one after};
\node[numbox,anchor=west] (N8) at (6.7,0 |- S8) {40 of 63 in-time collisions visible ($\ge4$ charged, $p_T>0.5$, $|\eta|<2.5$); 3 within 1 mm of the signal; the flat profile reweighted to the 2025 profile: $w=0.65$ at $\mu_{\rm true}=60$};
\node[numbox,anchor=west] (N9) at (6.7,0 |- S9) {\file{event_PU.lhe}: the chapter-7 rows + 961 vertex rows + 13\,456 particle rows + 175 jet rows + 651 constituent rows (1.7 MB)};
\end{tikzpicture}
}
\caption[The pileup stage as a flow chart, with the numbers of our
event]{The pileup stage as a flow chart. Left, the nine steps of
\file{standalone_pileup.py} with the names of the classes and functions
that do them; right, the numbers of our event at each step, from the
\code{-v} log of seed~1 at $\mup=60$ in $\bx=-12\ldots+3$.}
\label{fig:crowd_on_the_road:flow}
\end{figure}

\section{The energy ledger: our jet at this stage}
\label{sec:crowd_on_the_road:our_jet_at_this_stage_how_much_pileup_en}

The crowd has arrived, and the ledger of this stop has three lines where
the previous stops had one: what the photograph held, what the circle
reads with the crowd in it, and what is left after the subtraction.
\label{sec:crowd_on_the_road:record_format}
The record, \file{example/PU/event_PU.lhe}, keeps every row of
\cref{sec:family_photograph:record_format}, the listing, the vertex table
and the two jet tables of \cref{ch:family_photograph}, and adds four
tables after them, which the companion prints with a guide to their
columns. The \emph{summary} holds the scenario, the draws, the widths of
the beam spot and \rhoPU. The \emph{vertex listing} has one row per
collision of the sixteen crossings, $960$ of them and the signal first:
the crossing, the library event it came from, its process type, its
particles and its charged ones in the tracker's reach, whether the
tracker can see it, and its vertex, $x$ and $y$ in microns, $z$ in
centimetres and $t$ in nanoseconds with the $25\ns\times\bx$ offset
added; it is the \code{PileupSummaryInfo} and the vertex table of the
crowd in one. The \emph{particle listing} has every final particle of the
$63$ in-time collisions, $12\,113$ rows with the collision, the identity,
the four-momentum and the mixed jet that holds the particle, if any; the
out-of-time collisions are reproducible from their library events and are
not listed, since nothing at this stop uses their particles. The
\emph{jet listing} has one row per mixed \antikt jet, $173$ of them: its
momentum, its active area, $\rhoPU A$, the subtracted momentum, the
momentum of its signal and of its pileup constituents, the numbers of
each and of the collisions they came from, and the jet of
\cref{ch:family_photograph} it matches within $\dR<0.2$; and the
\emph{signal constituent listing} says, for each of the $651$ signal
particles that entered the clustering, which mixed jet holds it.

\begin{table}[htbp]\centering\small
\caption[Our jet and the $\bar b$ jet before and after the crowd]{Our jet
and the $\bar b$ jet before and after the crowd, from
\file{event_PU.lhe} as printed by \file{standalone_pileup.py}. The
photograph is the \akfour jet of \cref{ch:family_photograph}; the mixed
jet is the \antikt jet of the same $R$ on the signal and the in-time
crowd, matched to it within $\dR<0.2$; ``signal inside'' is the scalar
sum of the \pt of the family's particles in the mixed jet, the truth
the subtraction is judged against; the ensemble lines are the four
hundred draws of the crowd on the circle of $R=0.4$ around the
photograph's axis, for which $A=\pi R^2$.}
\label{tab:crowd_on_the_road:kinematics}
\scriptsize\setlength{\tabcolsep}{3pt}
\resizebox{\textwidth}{!}{%
\begin{tabular}{llrrrrrrr}\toprule
 & & \pt [GeV] & $m$ [GeV] & $n$ & $A$ & $\rhoPU A$ & $\pt-\rhoPU A$ & residual \\ \midrule
$b$ & the photograph, \akfour jet [0] (chapter 7) & $99.83$ & $14.83$ & $31$ & & & & \\
    & with the crowd: mixed jet [0] & $132.75$ & $45.27$ & $112$ ($31+81$) & $0.529$ & $36.09$ & $96.66$ & $-3.41$ \\
    & \quad signal inside / pileup inside & $100.07$ / $33.61$ & & $31$ / $81$ (34 coll.) & & & & \\
    & \quad the ensemble: the circle, 400 draws & $100.09$ + $27.8\pm8.4$ & & $65$ & $0.503$ & $29.7\pm5.8$ & & $-1.9\pm6.7$ \\ \midrule
$\bar b$ & the photograph, \akfour jet [1] & $94.09$ & $17.83$ & $18$ & & & & \\
    & with the crowd: mixed jet [1] & $124.71$ & $44.33$ & $98$ ($18+80$) & $0.499$ & $34.05$ & $90.66$ & $-3.64$ \\
    & \quad signal inside / pileup inside & $94.30$ / $31.26$ & & $18$ / $80$ (36 coll.) & & & & \\
    & \quad the ensemble: the circle, 400 draws & $93.2$ + $28.2\pm7.7$ & & & $0.503$ & $29.7\pm5.8$ & & $-1.5\pm5.9$ \\ \midrule
\multicolumn{2}{l}{the crossing: mixed jets above $3$ / $30\GeV$, after the subtraction} & \multicolumn{7}{l}{$173$ / $62$ / $6$; of the $62$, six are the family's and $56$ the crowd's own, none of which survives the subtraction}\\
\bottomrule
\end{tabular}}
\end{table}

\begin{ledger}{The Crowd on the Road}
  \ledgerrow{the photograph, \akfour jet [0] (chapter 7)}{99.83}{14.83}{31}{--}
  \ledgerrow{with the crowd: mixed jet [0], $\mup=60$}{132.75}{45.27}{112}{$+32.92$}
  \ledgerrow{after the subtraction, $\pt-\rhoPU A$}{96.66}{--}{112}{$-36.09$}
\end{ledger}

\Cref{tab:crowd_on_the_road:kinematics}, the ledger and
\cref{fig:crowd_on_the_road:journey} say what the record means, and the
lesson has two halves. The first is how much the crowd is: the circle
that held $99.83\GeV$ holds $132.75$, the $31$ particles of the family
have $81$ strangers among them from more than half the collisions of the
crossing, and the mass of the jet, which was $14.8\GeV$, is $45$, because
eighty soft particles spread over a circle of $0.4$ carry a great deal of
invariant mass and none of it is the family's; the $\bar b$ jet tells the
same story at $124.71$. The axis moved by $0.04$, the family stayed
whole, not one of its particles fell out and none joined, and that is
what the softness of the crowd means: it adds, it does not rearrange. In
\cref{fig:crowd_on_the_road:journey} the whole plane is khaki, the $173$
mixed jets tile it, $62$ of them above $30\GeV$ where the photograph had
six, and the family's circle is one circle among many until the
subtraction, after which the six of the family remain and not one of the
crowd's own: at $\mup=60$ a pileup jet above $30\GeV$ after $\rhoPU A$ is
a rare fluctuation, and the pileup jets that \cref{ch:the_police} will
have to identify live below that threshold. The second half is what the
subtraction does and does not do, and the ensemble lines of the table
carry it: the average is removed and two GeV of the family's own
underlying event with it, a bias the corrections of
\cref{ch:paying_for_uninvited} calibrate away, and seven GeV of
fluctuation remain and are not. The $\bar b$ jet at $\eta=-3.5$, in the
forward calorimeter where the crowd's plateau has begun to fall, reads
the same residual within a GeV, because at $\mup=60$ the $\eta$
dependence of the density is still smaller than its fluctuation; the
$\rhoPU(\eta)$ of \cite{Cacciari:2009dp}, or the $\eta$-dependent offset
of \cref{ch:paying_for_uninvited}, is the repair where it is not. The
crowd of \cref{fig:crowd_on_the_road:journey} is what the detector of
\cref{ch:crossing_the_border} will see, hits and all; the figure beside
it, the same frame with the crowd removed, is what it will try to get
back.

\begin{figure}[htbp]\centering
\includegraphics[width=0.495\textwidth]{PU/journey_PU_signal.pdf}\hfill
\includegraphics[width=0.495\textwidth]{PU/journey_PU.pdf}
\caption[The event before and with the crowd in the $(\eta,\phi)$
plane]{The event in the $(\eta,\phi)$ plane, in the layout of the
\CMSSW analyser pages. Left: the family alone, the photograph of
\cref{ch:family_photograph} drawn with the code of this stop. Right: the
same frame with the $63$ in-time collisions of our crossing, the crowd
in khaki, the signal particles coloured by the mixed jet that holds them,
the circles of the mixed jets with their momentum before and after the
subtraction, and our jet's line in the header: $99.83\to132.75\to96.66$.
Drawn by \file{plot_journey_PU.py} from \file{event_PU.lhe}.}
\label{fig:crowd_on_the_road:journey}
\end{figure}

\subsection{What this stop is, and what it is not}
\label{sec:crowd_on_the_road:caveats}

The agreement in \cref{tab:crowd_on_the_road:numbers} with \FASTJET is
exact where the input is identical, and the parts of the chapter that are
a model rather than a port carry their caveats:
\begin{itemize}
  \item \textbf{The crowd is \PYTHIA's.} The library is CP5's minimum
    bias, and the $0.95\GeV$ per unit area per collision is its number;
    a different tune moves it by ten percent, and the measured underlying
    event of \cite{CMS:2019csb} is what CP5 was tuned to. \rhoPU itself is
    measured in every event and does not care.
  \item \textbf{A generator-level crowd.} Every stable particle of every
    collision within $|\eta|<5$ counts here with its full momentum; the
    detector of \cref{ch:crossing_the_border} will lose the neutrals
    below its thresholds, the charged ones below $\pt\simeq0.3\GeV$ that
    curl in the field, and the forward ones in the cracks, so that the
    $\rhoPU$ of \cref{ch:paying_for_uninvited} is nearer $0.6\GeV$ per
    collision than $0.9$. The crowd shrinks; its structure does not.
  \item \textbf{One value of \mup, not the profile.} The record mixes at
    a fixed $\mup=60$, the \code{AVE\_60} of the companion's exercise; the
    production draws \mutrue flat between $55$ and $75$ and reweights,
    and a sample like this record cannot be reweighted at all. The
    ensemble draws the crowd at the same mean four hundred times, not over
    a year.
  \item \textbf{The out-of-time crowd is counted, not felt.} The $897$
    collisions of the other fifteen crossings are in the vertex listing
    with their offsets and nowhere else, because nothing at the generator
    level lasts $25\ns$. The calorimeters do, and \cref{ch:first_checkpoint}
    feels them.
  \item \textbf{The beam spot is a payload, not a model of the beams.}
    \Cref{eq:crowd_on_the_road:region} is what two head-on Gaussian
    bunches give, and \CMSSW's generator draws exactly that shape, but with
    the measured widths of early 2024 rather than widths computed from
    the bunches; a fill at the end of the year, with a different
    $\beta^*$ and crossing angle, had a different beam spot, and the real
    one is measured run by run.
  \item \textbf{The visible vertex is a proxy.} Four charged particles of
    $\pt>0.5\GeV$ within $|\eta|<2.5$ is a count, not a fit; the vertex
    finder of \cref{ch:first_checkpoint} needs tracks, and loses more of
    the crowd than the proxy does.
  \item \textbf{The subtraction is scalar.} $\pt-\rhoPU A$ corrects the
    transverse momentum and leaves the mass of the jet at its inflated
    value; the four-vector subtraction of \cite{Cacciari:2008gn}, with the
    area as a four-vector, repairs the mass too, and the substructure
    chapters will need it.
  \item \textbf{The ensemble draws the crowd, not the family.} The four
    hundred draws keep the photograph fixed and count the pileup inside a
    fixed circle, so that the $\pm7\GeV$ is the crowd's fluctuation
    alone; the mixed jet of the record, whose axis and area follow the
    crowd, adds a little to it.
  \item \textbf{One event is one sample.} The residual of $-3.41\GeV$ is
    this crossing's and the ensemble's is $-1.9\pm6.7$; the $63$
    collisions, the $34$ that reached our jet, the $68\GeV$ of \rhoPU
    against $59$ on average and the crowd vertex a twentieth of a
    millimetre from ours are this draw's.
\end{itemize}

\section{Further reading}
\label{sec:crowd_on_the_road:further_reading}
The jet area and the median density are \cite{Cacciari:2007fd}, with the
areas themselves, passive and active, in \cite{Cacciari:2008gn} and their
implementation in the \FASTJET manual~\cite{Cacciari:2011ma}; the
$\eta$-dependence of the density is \cite{Cacciari:2009dp}. The
inelastic cross section is measured in \cite{ATLAS:2016ygv,CMS:2018mlc}.
The simulation side is the \CMSSW premixing of \cite{Hildreth:2017caz}
and, for the experiment's own account of its crowd, the pileup
mitigation paper of CMS~\cite{CMS:2020ebo}, which measures \rhoPU and the
offset in Run-2 data. What comes next is in the particle-level methods,
\PUPPI~\cite{Bertolini:2014bba}, SoftKiller~\cite{Cacciari:2014gra},
constituent subtraction~\cite{Berta:2014eza} and jet
cleansing~\cite{Krohn:2013lba}, reviewed with a theorist's eye in
\cite{Soyez:2018opl}. The generator's side is \PYTHIA~8.3's
guide~\cite{Bierlich:2022pfr} for the soft processes and the tune
\cite{CMS:2019csb}. The accelerator of the open problems is the HL-LHC
design report~\cite{Apollinari:2017lhl}, with the crab cavities, the
levelling and the luminous region, and the detectors that will sort its
crowd are the Phase-2 tracker~\cite{CMS:2017lum}, the endcap
calorimeter~\cite{CMS:2017jpq} and the timing layer~\cite{Butler:2019rpu}.

\begin{pitfalls}
\begin{itemize}[nosep]
  \item \textbf{\mup is not the number of vertices.} $\mup=60$ is the
    mean of a Poisson draw; a given crossing has $60\pm8$ collisions, the
    tracker sees about $0.7\mup$ of them as vertices, and the
    \code{PileupSummaryInfo} holds all three numbers: the mean
    (\code{getTrueNumInteractions}), the draw
    (\code{getPU\_NumInteractions}) and, after reconstruction, $N_{\mathrm{PV}}$.
    Reweight on the mean, cut on the vertices, and never mix them.
  \item \textbf{40 MHz is not the crossing rate.} The bunch spacing is
    $25\ns$, but $2350$ of $3564$ slots collided in 2024: \cref{eq:crowd_on_the_road:mu}
    divides by $n_b f_{\mathrm{rev}}$, a third less, and $n_b$ changes
    with the filling scheme from year to year.
  \item \textbf{The library must match the signal.} The same beam spot,
    the same geometry, the same \GEANT, the same energy, and \emph{different
    seeds}: a library made with the default seeds of \code{cmsDriver.py}
    draws the same vertex sequence as the signal, and the first pileup
    event lands on the signal vertex in space and time. A library made
    with the HL-LHC beam spot under a 2024 signal would be wrong, and
    nothing in the job would notice.
  \item \textbf{The flat profile is not the data.} A sample mixed with
    \code{Run3\_Flat55To75} is meant to be reweighted to the measured
    profile of the year, on \mutrue; an analysis that forgets the weights
    compares data at $\langle\mup\rangle\simeq60$ with a simulation at
    $65$, and its jets are too wide and too many.
  \item \textbf{The median, not the mean.} A mean density over the plane
    includes the jets; the median does not. Conversely a median over too
    few patches, or over a region with more jets than empty patches, is a
    jet measurement in disguise.
  \item \textbf{\rhoPU is one number for the whole plane.} The crowd's
    plateau falls off at large $|\eta|$ and the detector's response with
    it; a jet in the endcap can be over-subtracted by a \rhoPU measured
    over $|\eta|<5$, and the offset corrections of
    \cref{ch:paying_for_uninvited} are binned in $\eta$ for that reason.
    \rhoPU also counts the signal's own underlying event, $3\GeV$ per
    unit area here, and subtracts it with the crowd.
  \item \textbf{$\pt-\rhoPU A$ can be negative.} A soft jet in a crowded
    event is subtracted to nothing or below; \FASTJET keeps the sign,
    \CMSSW floors the correction, and a threshold on the subtracted
    momentum is where the pileup jets of \cref{ch:the_police} start.
  \item \textbf{The area is not $\pi R^2$.} It is on average, for a lone
    hard jet; a neighbour within $2R$ takes a bite, a soft jet is irregular,
    and a jet's area fluctuates with the ghosts by a percent or two.
    Always use the active area the clustering measured.
  \item \textbf{Early and late are not in-time.} An out-of-time collision
    sits on the same beam spot $25\ns$ away; the tracker cannot tell,
    the calorimeter pulse integrates it, and a track and a cluster that
    disagree may be a crossing apart. The production window reaches
    twelve crossings back for that reason.
\end{itemize}
\end{pitfalls}

\begin{exercises}
\item Run the standalone script with \code{--mu 40} and with
  \code{--mu 75}, the edges of the Run-3 years, and compare \rhoPU, the
  pileup inside our jet and the number of pileup jets above $30\GeV$
  before the subtraction with the $\mup=60$ record. Is the pileup inside
  the circle linear in \mup? Is its fluctuation?
\item From \cref{eq:crowd_on_the_road:mu}, compute \mup for the LHC of
  2012 ($L=7.7\times10^{33}$, $n_b=1380$, $50\ns$ spacing,
  $\sigmainel=73\mbarn$ at $8\TeV$) and for 2018
  ($1.5\times10^{34}$, $2544$ bunches, $80\mbarn$), and compare with the
  $\sim21$ and $\sim37$ CMS recorded. What changes in the arithmetic when
  the spacing doubles?
\item The library has $2000$ events and a crossing uses $60$. Estimate
  the probability that a library event appears twice in the same crossing
  when read at random with replacement, and check it in the vertex listing
  of the record. Why does the production not worry about it?
\item Run the script with \code{--window 1 --sequential}, the conditions
  of the \CMSSW exercise of the companion, and compare the three draws,
  the particles per crossing and the vertex spread with the companion's
  table of mixed events. Which numbers can agree, and which cannot?
\item With \code{fill\_profile} and \code{year\_profile}, build a profile
  of a year whose fills are levelled at $62$ and compute the weights of
  \cref{eq:crowd_on_the_road:reweight} for the flat $55$--$75$ sample bin
  by bin. What fraction of the simulated events carries a weight below
  $0.1$, and what does that cost in statistical power?
\item The time spread of the beam spot is $120\ps$. Treating the
  collisions as Gaussian in time, what fraction of the crowd is within
  $\pm30\ps$ of the signal, on average? Check it against the
  \code{within 30 ps} count of the ensemble, and say what a timing layer
  with that resolution would have left of the in-time crowd.
\item Compute the grid median of the signal alone with
  \code{grid\_rho} and explain why it is $3.4\GeV$ per unit area and not
  zero. What would the $k_T$-jet median give on the signal alone, and why?
\item Draw the pileup inside the circle against the number of in-time
  collisions over the ensemble and fit a line. Compare the slope with
  $\rhoPU_1\pi R^2$ from \cref{eq:crowd_on_the_road:rho1}, and the
  scatter around the line with $\sqrt{n}$ times the per-collision
  fluctuation.
\item Replace the grid median by the mean over the cells and redo the
  subtraction of the twelve leading jets of the record. Which jets get
  worse, and by how much?
\item Change \code{GHOST\_AREA} to $0.04$ and to $0.0025$ and compare the
  areas of the leading jets and the time of the clustering. At what ghost
  area does the area of our jet stop changing within its own fluctuation?
\item A jet of $R=0.8$ has four times the area. Using the ensemble's
  numbers, estimate the pileup inside an $R=0.8$ circle around our jet and
  its fluctuation, and say why the boosted-top jets of
  \cref{ch:three_families} cannot live with $\rhoPU A$ alone.
\item Run \code{--scenario ultimate} and compare the four numbers of
  this chapter, \rhoPU, the pileup inside the circle, the residual and
  the number of pileup jets above $30\GeV$ after the subtraction, between
  $\mup=60$ and $200$. Which of them scale with \mup, which with
  $\sqrt{\mup}$, and which do neither?
\end{exercises}

\subsection*{Open problems in the context of Phase-2 LHC}
\label{sec:crowd_on_the_road:open_problems}

The crowd of this chapter is the crowd of 2024; the HL-LHC will bring
three times more of it, two hundred collisions per crossing at $14\TeV$,
and the detector CMS is building for it, with a tracker to $|\eta|=4$, a
high-granularity endcap calorimeter and a timing layer, is designed
around that number. The scripts of this stop keep that road one option
away, \code{--scenario ultimate}: the same family, $\mup=200$ in
$\bx=-3\ldots+3$, the luminous region of the crab cavities, and the
questions below can each be started with it.
\begin{enumerate}
  \item \textbf{The resolution floor.} The $\pm7\GeV$ left on a
    $100\GeV$ jet after $\rhoPU A$ at $\mup=60$ becomes $\pm14$ at
    $\mup=200$ at the generator level, and the real floor is what the jet
    energy resolution of \cite{Azzi:2019yne} assumes for every HL-LHC
    measurement with jets. With the ensemble at $\mup=60$, $140$ and
    $200$ and the per-collision density of
    \cref{eq:crowd_on_the_road:rho1}, decompose the fluctuation into its
    Poisson part, the multiplicity part and the geometric part, and ask
    which of the three the tracks of \cref{ch:the_police} and the
    $30\ps$ of the MTD~\cite{Butler:2019rpu} can remove; the open
    question is the size of the irreducible neutral part at $|\eta|>2.5$,
    where there are no tracks and the HGCAL~\cite{CMS:2017jpq} has only
    its own timing.
  \item \textbf{The density in $\eta$ and in time.} \rhoPU is one number
    here and one number in \CMSSW; the crowd's plateau is not flat and
    the detector's response to it less so, and at $\mup=200$ the
    $\eta$-dependent offset of \cref{ch:paying_for_uninvited} becomes
    the largest correction of the chain. With the library, measure
    $\rhoPU_1(\eta)$ and its fluctuation in bins of $\eta$, and ask
    whether a $\rhoPU(\eta,t)$, a density per unit area per unit time
    that the timing detectors could measure, subtracts better than
    $\rhoPU(\eta)$ for a jet whose own time is known; the question, raised
    for the MTD in \cite{Butler:2019rpu}, is open for the calorimeter jets
    that carry no track time.
  \item \textbf{Pileup jets.} At $\mup=60$ not one of the $56$ crowd jets
    above $30\GeV$ survives the subtraction; at $\mup=200$ a dozen do,
    twice as many as the family's, and every analysis with a jet veto or
    a jet count will live on the identification of
    \cref{ch:the_police}. With the jet listings of the two records,
    characterise the survivors, their number of collisions, their widths
    and their charged fractions, and ask which variable separates a crowd
    jet from a soft quark jet at the generator level; the Phase-2 question
    is whether the L1 track trigger of \cite{CMS:2020jkc} can make that
    separation at $40\,\mathrm{MHz}$.
  \item \textbf{The luminous region and the crab cavities.} The beam
    spot of 2024 is a payload fitted to data; the HL-LHC region is a
    model built from the bunches, in which the crossing angle would halve
    the overlap and the crab cavities give it back, and no data exist yet
    to fit. With \code{LuminousRegion.density} and its switches, compute
    the luminosity lost without the crab cavities and the length and the
    time spread of the region as a function of the crab angle, and ask how
    the vertex finder and the MTD should be tuned if the crabbing scheme
    changes during a fill, as the design report~\cite{Apollinari:2017lhl}
    foresees.
  \item \textbf{The library, the premixing and the profile at
    $\mup=200$.} The production premixes because the classical mixing
    cannot read $1400$ events per signal event, and the premixed sample
    loses the hit-level truth of the crowd that this chapter's pictures
    rely on; the HL-LHC profile is a spike at the levelled value with a
    tail, and a flat band like Run~3's would waste most of its events.
    With the standalone mixing, which keeps the truth at no cost, measure
    how many distinct library events a crossing needs before the reuse of
    a collision within one event becomes visible in the pileup-jet rate,
    and with \code{fill\_profile} build the mixing histogram that
    minimises the statistical cost of the reweighting for a levelled
    fill; the open question, for the Phase-2 production of
    \cite{Hildreth:2017caz}, is how small a premixed library may be
    before correlations between signal events appear in a measurement.
\end{enumerate}
The next stop takes the family and the crowd across the border, into the
detector: the magnetic field, the material, the hits, and the first
things the crowd does to them.
